% Standalone research note for the hybrid-local-solver project.

\documentclass[11pt]{article}

% Common class-neutral preamble for standalone notes under manuscript/notes/.
% Note entry points all live exactly two directories below manuscript/.

\usepackage[margin=1in]{geometry}
\usepackage{amsmath,amssymb,amsthm,mathtools,bm}
\usepackage{algorithm}
\usepackage{algorithmic}
\usepackage{booktabs,tabularx,multirow,array,longtable}
\usepackage{enumitem}
\usepackage{microtype}
\usepackage[numbers,sort&compress]{natbib}
\usepackage{xcolor}
\usepackage[colorlinks=true,citecolor=blue,linkcolor=black,urlcolor=blue]{hyperref}
\usepackage[capitalize,nameinlink,noabbrev]{cleveref}

\newtheorem{theorem}{Theorem}[section]
\newtheorem{lemma}[theorem]{Lemma}
\newtheorem{proposition}[theorem]{Proposition}
\newtheorem{corollary}[theorem]{Corollary}
\newtheorem{conjecture}[theorem]{Conjecture}
\theoremstyle{definition}
\newtheorem{definition}[theorem]{Definition}
\newtheorem{assumption}[theorem]{Assumption}
\newtheorem{problem}[theorem]{Problem}
\newtheorem{observation}[theorem]{Observation}
\theoremstyle{remark}
\newtheorem{remark}[theorem]{Remark}
\theoremstyle{plain}

% Canonical mathematical notation for the active manuscript.
%
% This file preserves the recurring notation style used in the author's
% NeurIPS 2024 and 2025 papers while excluding unrelated tensor and
% machine-learning boilerplate from those archived command collections.
%
% Prefix convention:
%   \v...  bold vectors
%   \m...  bold matrices
%   \g...  calligraphic sets, graphs, and families
%   \s...  blackboard-bold spaces and number systems

\newcommand{\method}{\textsc{HybridLocalSolver}}

% Font and scalar shorthands retained from the archived manuscripts.
\newcommand{\mc}{\mathcal}
\newcommand{\R}{\mathbb{R}}
\newcommand{\E}{\mathbb{E}}
\newcommand{\G}{\mathbb{G}}
\newcommand{\N}{\mathcal{N}}
% Accuracy namespaces are semantic: do not add a bare \eps alias.
% Degree-normalized residual tolerance of classical APPR is kept distinct from
% the objective-gap and proximal fixed-point tolerances of the RPPR sections.
\newcommand{\epsappr}{\varepsilon_{\mathrm{appr}}}
\newcommand{\epsobj}{\varepsilon_{\mathrm{obj}}}
\newcommand{\epspg}{\varepsilon_{\mathrm{pg}}}
\newcommand{\epsppr}{\varepsilon_{\mathrm{ppr}}}
\newcommand{\epsin}{\varepsilon_{\mathrm{in}}}
\newcommand{\epskkt}{\varepsilon_{\mathrm{kkt}}}
\newcommand{\epsburn}{\varepsilon_{\mathrm{burn}}}
\newcommand{\epssol}{\varepsilon_{\mathrm{sol}}}
\newcommand{\epspprinst}[1]{\varepsilon_{\mathrm{ppr},#1}}
\newcommand{\trans}{\mathsf{T}}
\newcommand{\grad}{\nabla}

% Delimiters and generic helpers.
% Norms and inner products use fixed-size delimiters by default.
% Use an optional size (e.g. \norm[\big]{...}) or * for automatic sizing.
\DeclarePairedDelimiter{\norm}{\lVert}{\rVert}
\newcommand{\abs}[1]{\left\lvert #1 \right\rvert}
\DeclarePairedDelimiterX{\inner}[2]{\langle}{\rangle}{#1, #2}
\newcommand{\ip}[2]{\inner{#1}{#2}}
\newcommand{\card}[1]{\left\lvert #1 \right\rvert}
\newcommand{\set}[1]{\left\{#1\right\}}
\newcommand{\ceil}[1]{\left\lceil #1 \right\rceil}
\newcommand{\floor}[1]{\left\lfloor #1 \right\rfloor}
\newcommand{\comp}[1]{\overline{#1}}
\newcommand{\conj}[1]{\overline{#1}}
\newcommand{\mean}[1]{\overline{#1}}
\newcommand{\bigO}{\mathcal{O}}
\newcommand{\softO}{\widetilde{\mathcal{O}}}
\newcommand{\softOmega}{\widetilde{\Omega}}
\newcommand{\pos}[1]{\left[#1\right]_{+}}

% Bold lowercase vectors.
\newcommand{\va}{\bm{a}}
\newcommand{\vb}{\bm{b}}
\newcommand{\vc}{\bm{c}}
\newcommand{\vd}{\bm{d}}
\newcommand{\ve}{\bm{e}}
\newcommand{\vf}{\bm{f}}
\newcommand{\vg}{\bm{g}}
\newcommand{\vh}{\bm{h}}
\newcommand{\vi}{\bm{i}}
\newcommand{\vj}{\bm{j}}
\newcommand{\vk}{\bm{k}}
\newcommand{\vl}{\bm{l}}
\newcommand{\vm}{\bm{m}}
\newcommand{\vn}{\bm{n}}
\newcommand{\vo}{\bm{o}}
\newcommand{\vp}{\bm{p}}
\newcommand{\vq}{\bm{q}}
\newcommand{\vr}{\bm{r}}
\newcommand{\vs}{\bm{s}}
\newcommand{\vt}{\bm{t}}
\newcommand{\vu}{\bm{u}}
\newcommand{\vv}{\bm{v}}
\newcommand{\vw}{\bm{w}}
\newcommand{\vx}{\bm{x}}
\newcommand{\vy}{\bm{y}}
\newcommand{\vz}{\bm{z}}

% Bold Greek vectors and special vectors.
\newcommand{\valpha}{\bm{\alpha}}
\newcommand{\vbeta}{\bm{\beta}}
\newcommand{\vdelta}{\bm{\delta}}
\newcommand{\vepsilon}{\bm{\epsilon}}
\newcommand{\veta}{\bm{\eta}}
\newcommand{\vgamma}{\bm{\gamma}}
\newcommand{\vlambda}{\bm{\lambda}}
\newcommand{\vmu}{\bm{\mu}}
\newcommand{\vomega}{\bm{\omega}}
\newcommand{\vphi}{\bm{\phi}}
\newcommand{\vpi}{\bm{\pi}}
\newcommand{\vpsi}{\bm{\psi}}
\newcommand{\vrho}{\bm{\rho}}
\newcommand{\vsigma}{\bm{\sigma}}
\newcommand{\vtheta}{\bm{\theta}}
\newcommand{\vxi}{\bm{\xi}}
\newcommand{\vell}{\bm{\ell}}
\newcommand{\vDelta}{\bm{\Delta}}
\newcommand{\vGamma}{\bm{\Gamma}}
\newcommand{\vLambda}{\bm{\Lambda}}
\newcommand{\vPhi}{\bm{\Phi}}
\newcommand{\vSigma}{\bm{\Sigma}}
\newcommand{\vzero}{\bm{0}}
\newcommand{\vone}{\bm{1}}
\newcommand{\one}{\mathbf{1}}
\newcommand{\zero}{\vzero}
\newcommand{\eunit}[1]{\bm{e}_{#1}}
\newcommand{\vDeltax}{\bm{\Delta}\bm{x}}

% Bold uppercase matrices.
\newcommand{\mA}{\bm{A}}
\newcommand{\mB}{\bm{B}}
\newcommand{\mC}{\bm{C}}
\newcommand{\mD}{\bm{D}}
\newcommand{\mE}{\bm{E}}
\newcommand{\mF}{\bm{F}}
\newcommand{\mG}{\bm{G}}
\newcommand{\mH}{\bm{H}}
\newcommand{\mI}{\bm{I}}
\newcommand{\mJ}{\bm{J}}
\newcommand{\mK}{\bm{K}}
\newcommand{\mL}{\bm{L}}
\newcommand{\mM}{\bm{M}}
\newcommand{\mN}{\bm{N}}
\newcommand{\mO}{\bm{O}}
\newcommand{\mP}{\bm{P}}
\newcommand{\mQ}{\bm{Q}}
\newcommand{\mR}{\bm{R}}
\newcommand{\mS}{\bm{S}}
\newcommand{\mT}{\bm{T}}
\newcommand{\mU}{\bm{U}}
\newcommand{\mV}{\bm{V}}
\newcommand{\mW}{\bm{W}}
\newcommand{\mX}{\bm{X}}
\newcommand{\mY}{\bm{Y}}
\newcommand{\mZ}{\bm{Z}}

% Bold Greek matrices retained from the graph papers.
\newcommand{\mBeta}{\bm{\beta}}
\newcommand{\mDelta}{\bm{\Delta}}
\newcommand{\mGamma}{\bm{\Gamma}}
\newcommand{\mLambda}{\bm{\Lambda}}
\newcommand{\mOmega}{\bm{\Omega}}
\newcommand{\mPhi}{\bm{\Phi}}
\newcommand{\mPi}{\bm{\Pi}}
\newcommand{\mSigma}{\bm{\Sigma}}
\newcommand{\mTheta}{\bm{\Theta}}

% Calligraphic symbols.
\newcommand{\gA}{\mathcal{A}}
\newcommand{\gB}{\mathcal{B}}
\newcommand{\gC}{\mathcal{C}}
\newcommand{\gD}{\mathcal{D}}
\newcommand{\gE}{\mathcal{E}}
\newcommand{\gF}{\mathcal{F}}
\newcommand{\gG}{\mathcal{G}}
\newcommand{\gH}{\mathcal{H}}
\newcommand{\gI}{\mathcal{I}}
\newcommand{\gJ}{\mathcal{J}}
\newcommand{\gK}{\mathcal{K}}
\newcommand{\gL}{\mathcal{L}}
\newcommand{\gM}{\mathcal{M}}
\newcommand{\gN}{\mathcal{N}}
\newcommand{\gO}{\mathcal{O}}
\newcommand{\gP}{\mathcal{P}}
\newcommand{\gQ}{\mathcal{Q}}
\newcommand{\gR}{\mathcal{R}}
\newcommand{\gS}{\mathcal{S}}
\newcommand{\gT}{\mathcal{T}}
\newcommand{\gU}{\mathcal{U}}
\newcommand{\gV}{\mathcal{V}}
\newcommand{\gW}{\mathcal{W}}
\newcommand{\gX}{\mathcal{X}}
\newcommand{\gY}{\mathcal{Y}}
\newcommand{\gZ}{\mathcal{Z}}

% Blackboard-bold symbols.
\newcommand{\sA}{\mathbb{A}}
\newcommand{\sB}{\mathbb{B}}
\newcommand{\sC}{\mathbb{C}}
\newcommand{\sD}{\mathbb{D}}
\newcommand{\sE}{\mathbb{E}}
\newcommand{\sF}{\mathbb{F}}
\newcommand{\sG}{\mathbb{G}}
\newcommand{\sH}{\mathbb{H}}
\newcommand{\sI}{\mathbb{I}}
\newcommand{\sJ}{\mathbb{J}}
\newcommand{\sK}{\mathbb{K}}
\newcommand{\sL}{\mathbb{L}}
\newcommand{\sM}{\mathbb{M}}
\newcommand{\sN}{\mathbb{N}}
\newcommand{\sO}{\mathbb{O}}
\newcommand{\sP}{\mathbb{P}}
\newcommand{\sQ}{\mathbb{Q}}
\newcommand{\sR}{\mathbb{R}}
\newcommand{\sS}{\mathbb{S}}
\newcommand{\sT}{\mathbb{T}}
\newcommand{\sU}{\mathbb{U}}
\newcommand{\sV}{\mathbb{V}}
\newcommand{\sW}{\mathbb{W}}
\newcommand{\sX}{\mathbb{X}}
\newcommand{\sY}{\mathbb{Y}}
\newcommand{\sZ}{\mathbb{Z}}

% Frequently used operators.
\DeclareMathOperator{\Cov}{Cov}
\DeclareMathOperator{\Var}{Var}
\DeclareMathOperator{\diag}{diag}
\DeclareMathOperator{\nei}{Nei}
\DeclareMathOperator{\nnz}{nnz}
\DeclareMathOperator{\pr}{pr}
\DeclareMathOperator{\prox}{prox}
\DeclareMathOperator{\push}{push}
\DeclareMathOperator{\res}{res}
\DeclareMathOperator{\rel}{Rel}
\DeclareMathOperator{\relu}{ReLU}
\DeclareMathOperator{\sign}{sign}
\DeclareMathOperator{\supp}{supp}
\DeclareMathOperator{\tr}{tr}
\DeclareMathOperator{\Tr}{Tr}
\DeclareMathOperator{\vol}{vol}
\DeclareMathOperator{\work}{work}
\DeclareMathOperator{\Work}{Work}
\DeclareMathOperator*{\argmax}{arg\,max}
\DeclareMathOperator*{\argmin}{arg\,min}

% Project-wide names and claim-status labels used by standalone research notes.
% Mathematical notation belongs in math_commands.tex.

% Algorithms and solver families.
\newcommand{\AESP}{\textsc{AESP}}
\newcommand{\APPR}{\textsc{APPR}}
\newcommand{\ASPR}{\textsc{ASPR}}
\newcommand{\FISTA}{\textsc{FISTA}}
\newcommand{\ISTA}{\textsc{ISTA}}
\newcommand{\LOCAPPR}{\textsc{LocAPPR}}
\newcommand{\LOCGD}{\textsc{LocGD}}
\newcommand{\LOCSOR}{\textsc{LocSOR}}
\newcommand{\LOCCH}{\textsc{LocCH}}
\newcommand{\LOCHB}{\textsc{LocHB}}
\newcommand{\LocProxGS}{\textsc{LocProxGS}}
\newcommand{\Hybrid}{\textsc{AESP--LocSOR}}

% Claim classifications required by the repository research-note protocol.
\newcommand{\statussource}{\textbf{Source result}}
\newcommand{\statusproved}{\textbf{Proved here}}
\newcommand{\statusconditional}{\textbf{Conditional}}
\newcommand{\statusconjecture}{\textbf{Open / conjectural}}
\newcommand{\statusopen}{\textbf{Open}}
\newcommand{\statusempirical}{\textbf{Empirical}}
\newcommand{\statusobservation}{\textbf{Empirical observation}}
\newcommand{\statusrefuted}{\textbf{Refuted route}}

% Compatibility names used by the older complete-note presentation layer.
\newcommand{\Status}[2]{\textbf{#1:} #2}
\newcommand{\SourceStatus}{\textnormal{\textsc{Source result}}}
\newcommand{\ProvedStatus}{\textnormal{\textsc{Proved here}}}
\newcommand{\ConditionalStatus}{\textnormal{\textsc{Conditional}}}
\newcommand{\ComputedStatus}{\textnormal{\textsc{Computed}}}
\newcommand{\RefutedStatus}{\textnormal{\textsc{Refuted route}}}
\newcommand{\OpenStatus}{\textnormal{\textsc{Open}}}


\usepackage{longtable}

\title{Response-Preconditioned Local PageRank:\
A Common Interface for Iterative Repair and Persistent Schur State}
\author{Baojian Zhou\\
\small Working research note for \texttt{hybrid-local-solver}}
\date{August 2026}

\begin{document}
\maketitle

\begin{abstract}
The project currently has two credible routes to faster local PageRank.  Local
first-order and Krylov methods seek work
$\widetilde O(V/\sqrt\alpha)$ on a certified active envelope of charged volume
$V$.  Growing-active-set SDD and directed-response methods can remove every
polynomial dependence on $1/\alpha$ on paths and several structured graph
classes, and could reach $\widetilde O(V)$ if their evolving boundary response
were maintainable at output-sensitive cost.  Neither endpoint is known on
arbitrary graphs.

This note formulates their common operator and a response-preconditioned
composition.  A settled anchor stores an implicit approximation to its
restricted inverse, newly admitted vertices form a frontier buffer, and
iterative repair acts on the preconditioned current face.  The proved
backend-independent facts include the exact bordered Schur correction, an
energy-to-coordinate boundary interval with graph-universal leverage at most
$\sqrt{(1-\alpha)/2}$, a source-aware leverage-packing theorem, a
Schur-diagonal-loss identity with a per-label temporal telescope, a
margin-adaptive hierarchical-locator reduction, a two-ledger square-root
wake-up bound, a grounded-Laplacian reduction to dynamic electrical-flow
location, exposed-incidence sufficiency, a transposed fixed-anchor harmonic
sketch identity, and geometric rebuild amortization.  Two
complementary response
ledgers sharpen the reporting target.  The first is a monotone-variation ledger:
inactive normalized boundary demands move in only one direction and have
total degree-weighted variation at most $(1-\alpha)/2$.  The second is an
energy-packing ledger: the squared exterior demand change is
at most the exact primal shock.  Under an exhaustively charged response and
boundary interface, these yield a conditional work
theorem of the form
\[
 \widetilde O\bigl(
 U_{\rm resp}(V)+V\sqrt{\kappa_{\rm eff}}+V
 \bigr).
\]
The identity preconditioner recovers the iterative product target; a
near-linear response interface with bounded effective condition approaches the
output scale.  The theorem is not a graph-uniform construction.  The missing
ingredient remains certified finite-band response maintenance on general
nonequitable cyclic cores, together with a no-restart expanding-subspace
accelerated repair if the response is inexact.
\end{abstract}

\tableofcontents

\section{Scope and claim ledger}
\label{sec:scope}

% Canonical source-aligned PageRank/RPPR reference formulation.
%
% This fragment is intentionally heading-free at the section level so that the
% active paper and every standalone note can input it. It is a manuscript
% reference convention, not an implementation-wide residual decision.
% Primary source: Fountoulakis and Martinez-Rubio, arXiv:2602.21138v2,
% PDF p. 3, Section 3 and equation (RPPR).

\subsection{Common source-aligned PageRank and RPPR model}
\label{subsec:shared-source-aligned-problem}

All documents in this manuscript workspace use the following reference
language. It follows the plain italic notation of the comparison paper:
\(x,s,A,D,Q,I\) denote column vectors or matrices as determined by their
displayed domains. This is the scoped exception to the project's usual
bold-vector and bold-matrix typography. A note may impose a stronger
assumption after this block, but it must not redefine these objects.

Let \(G=(V,E)\) be a finite, simple, undirected, unweighted graph without
isolated vertices, let \(V=[n]:=\{1,\ldots,n\}\), and let
\(A\in\{0,1\}^{n\times n}\) be its symmetric adjacency matrix. For
\(i\in V\), write
\[
    \mathcal N(i):=\{j\in V:\{i,j\}\in E\},
    \qquad
    d_i:=|\mathcal N(i)|,
    \qquad
    D:=\diag(d_1,\ldots,d_n).
\]
For \(S\subseteq V\), define
\[
    \mathcal N(S):=\bigcup_{i\in S}\mathcal N(i),
    \qquad
    \partial S:=\mathcal N(S)\setminus S,
    \qquad
    \operatorname{Ext}(S):=V\setminus(S\cup\partial S),
\]
\begin{equation}
    \vol(S):=\sum_{i\in S}d_i.
    \label{eq:shared-volume}
\end{equation}
The support of a vector is
\(\supp(x):=\{i\in[n]:x_i\neq0\}\). Matrix inequalities use the Loewner
order, and all unqualified vector inequalities are coordinatewise.

The seed is a column distribution
\begin{equation}
    s\in\R_+^n,
    \qquad
    \inner{\one}{s}=1.
    \label{eq:shared-seed}
\end{equation}
The single-seed specialization is always written \(s=e_v\), where \(v\in V\)
is the seed vertex. Thus \(s\) is never used as a vertex label. For
\(\alpha\in(0,1]\), define
\begin{equation}
\begin{aligned}
    \mathcal L&:=I-D^{-1/2}AD^{-1/2},\\
    Q&:=\alpha I+\frac{1-\alpha}{2}\mathcal L
      =\frac{1+\alpha}{2}I
       -\frac{1-\alpha}{2}D^{-1/2}AD^{-1/2},\\
    b&:=\alpha D^{-1/2}s.
\end{aligned}
    \label{eq:shared-pagerank-matrices}
\end{equation}
Since \(0\preceq\mathcal L\preceq2I\),
\(\alpha I\preceq Q\preceq I\). The shared smooth PageRank quadratic is
\begin{equation}
    f(x):=\frac12\inner{x}{Qx}-\inner{b}{x},
    \qquad
    \nabla f(x)=Qx-b.
    \label{eq:shared-pagerank-objective}
\end{equation}
Its unique unregularized minimizer and associated lazy personalized PageRank
vector are
\begin{equation}
    x^0:=Q^{-1}b,
    \qquad
    \pi:=D^{1/2}x^0.
    \label{eq:shared-ppr-solution}
\end{equation}

For a regularization parameter \(\rho>0\), reserve \(g_\rho\) for the
nonsmooth weighted \(\ell_1\) term and define
\begin{equation}
    g_\rho(x):=\alpha\rho\norm{D^{1/2}x}_1,
    \qquad
    F_\rho(x):=f(x)+g_\rho(x),
    \qquad
    x^\star(\rho):=\argmin_{x\in\R^n}F_\rho(x).
    \label{eq:shared-rppr-objective}
\end{equation}
When \(\rho\) is fixed, \(x^\star\) abbreviates \(x^\star(\rho)\), never the
unregularized point \(x^0\). Write
\[
    S^\star(\rho):=\supp(x^\star(\rho)),
    \qquad
    I^\star(\rho):=[n]\setminus S^\star(\rho).
\]
The source records \(x^\star(\rho)\geq0\),
\(\vol(S^\star(\rho))\leq1/\rho\), and the coordinatewise conditions
\begin{equation}
    \nabla_i f(x^\star(\rho))
    \in
    \begin{cases}
        \{-\alpha\rho\sqrt{d_i}\}, & x_i^\star(\rho)>0,\\
        [-\alpha\rho\sqrt{d_i},0], & x_i^\star(\rho)=0.
    \end{cases}
    \label{eq:shared-rppr-kkt}
\end{equation}

The common computational unit is one repeated adjacency-list scan: scanning
coordinate \(i\) costs \(d_i\), and scanning a set \(S\) costs \(\vol(S)\).
Iteration count and wall-clock time do not replace this degree-weighted work
measure.

Finally, accuracy symbols remain distinct. The APPR threshold is
\(\epsappr\); \(\epsobj\) denotes an objective-gap target;
\(\epspg\) denotes the source experiment's proximal fixed-point tolerance; and
\(\epsppr\) may denote a document-scoped degree-normalized PPR error target.
No equality or conversion among these quantities is assumed. In particular,
this shared model does not resolve the repository-wide residual decision.


The shared project model supplies the symmetric Stieltjes matrix $Q$, the
smooth PageRank quadratic, and the RPPR objective.  This note uses the shared
model only as a mathematical reference.  It does not adopt a repository-wide
residual or identify the ACL, objective-gap, proximal, and degree-normalized
accuracy namespaces.

{\small
\begin{longtable}{@{}p{0.19\textwidth}p{0.75\textwidth}@{}}
\toprule
Status & Statement \\
\midrule
\endfirsthead
\toprule
Status & Statement \\
\midrule
\endhead
\statusproved & Iterative and response methods act on the same restricted
Green operator $Q_{SS}^{-1}$.  A bordered active-set expansion has an exact
Schur correction. \\
\statusproved & Successive exact corrections on nested monotone faces are
$Q$-orthogonal.  A primal correction can be dense on the enlarged face, but
its dual image under $Q$ vanishes on the old face and is supported only on the
new batch and current boundary. \\
\statusproved & The entire lifted Schur-frontier spaces are mutually
$Q$-orthogonal.  Arbitrary certified frontier-solve errors therefore add in
quadrature, so the settled face needs no numerical restart when its response
is applied at charged cost. \\
\statusproved & Consecutive exact light transfers may be batched
associatively: their dual signatures add, their endpoint and exterior demands
are unchanged, and sequential Schur elimination equals one aggregate
elimination. \\
\statusproved & An energy-norm error certificate gives a rigorous interval for
every queried boundary KKT demand.  Every response leverage is at most
$\sqrt{(1-\alpha)/2}$, so coordinate-specific leverage estimation is
unnecessary. \\
\statusproved & If a correction is generated by a frontier source batch $B$,
its source-aware squared boundary leverages sum to at most $|B|$.  This gives
an a priori output-volume bound on ambiguous intervals without materializing
the dense primal correction. \\
\statusproved & Given constant-factor aggregate source-aware leverage masses
on a bounded-arity boundary hierarchy, branch and bound locates every
potentially ambiguous row in a number of group probes controlled by $|B|$ and
the correction energy, rather than the ambient boundary size. \\
\statusproved & Every source-aware squared leverage is exactly the decrease of
one terminal Schur diagonal across the batch elimination.  These losses
telescope to at most $(1-\alpha)/2$ per still-exterior label.  A one-sided
group oracle needs only additive accuracy at the current pruning scale, and
can replace the source-response sketches entirely. \\
\statusproved & Schur-diagonal loss and orthogonal correction energy can be
combined multiplicatively.  If a coordinate or hierarchy node is awakened
only when the product interval reaches its current margin, the sum of its
wake-up margins is bounded by the geometric mean of its total diagonal loss
and the global correction-energy budget.  For a unit seed, this produces a
per-coordinate $1/\sqrt\alpha$ rather than $1/\alpha$ scale at margin
$\Theta(\alpha\tau)$. \\
\statusproved & The shared PageRank matrix is diagonally congruent to the
grounded Laplacian of the original graph with conductance
$(1-\alpha)/2$ on each graph edge and conductance $\alpha d_v$ from $v$ to
ground.  Principal Schur complements and normalized exterior responses obey
the same congruence. \\
\statusproved & The response rows, source-aware leverages, and Schur-diagonal
losses for a current face depend only on that face's internal edges, its cut
incidences, and the degrees of the incident boundary labels.  Edges wholly in
the unexposed exterior are irrelevant.  Thus static locator rebuilding can be
charged to geometrically growing exposed volume; dynamic update and query
costs remain open. \\
\statusproved & At a fixed anchor, any linear sketch of the full normalized
exterior dual signature has an exact transposed harmonic-response formula.
After precomputing one anchor response per sketch row, a frontier solve updates
the sketch from frontier-incidence scans and stored anchor-sketch row reads,
without materializing the dense old-face correction or scanning the boundary.
\\
\statusproved & A logarithmic batch of source-normalized Gaussian response
sketches simultaneously supplies constant-factor group masses for every node
of a fixed boundary hierarchy.  Dynamic local maintenance of those sketches
remains open. \\
\statusproved & Fresh per-batch normalized response sketches can be accumulated
causally along an adaptive trace.  Their squared masses remain valid for every
prefix, and a newly exposed boundary label has zero contribution from all past
batches, so no response replay is required. \\
\statusproved & Across a nested exact-face expansion, the squared exterior
demand change is at most the primal energy shock.  Boundary degree volume
moving by normalized margin $\theta$ is therefore at most
$2\Delta/\theta^2$. \\
\statusproved & Along a monotone support-safe exact-face trace, every inactive
normalized boundary demand is nondecreasing and the total degree-weighted
variation is at most $(1-\alpha)/2$.  Exact threshold crossings therefore
occur at most once per boundary coordinate. \\
\statusproved & If anchor volumes grow geometrically, all linear-volume
rebuild costs sum to a constant multiple of the final volume. \\
\statusproved & A converged frontier probe may precede a heavy rebuild without
changing the eventual exact response.  It can omit a terminal rebuild when the
probe already certifies the requested output. \\
\statusconditional & Given a fully charged response interface and
preconditioned repairs on geometrically growing faces, total work is bounded
by response maintenance plus final volume times the square root of the
effective preconditioned condition number, up to logarithms. \\
\statusopen & No arbitrary-graph implementation of the required finite-band
heavy-change reporter is proved.  The remaining primitive is dynamic local
maintenance of additive aggregate terminal-Schur diagonal losses, or the
equivalent squared-response range-add, without visiting every boundary leaf.
\\
\statusrefuted & The phrase ``warm start the SDD solver'' is not a sufficient
algorithm: explicit old-face correction, complete boundary refresh, or full
vector output after every expansion can recreate quadratic repeated-prefix
work. \\
\bottomrule
\end{longtable}
}

The architectural recommendation is therefore risk-adjusted rather than
absolute.  A graph-uniform output-linear response solver would dominate the
mixed construction asymptotically.  Conversely, a graph-uniform
expanding-subspace first-order theorem could attain the product scale without
nontrivial response state.  The mixed route is useful because it exposes and
combines the partial theorems already available at both endpoints.

\section{The common restricted operator}
\label{sec:operator}

For a nonnegative RPPR face $S$, define the note-scoped shifted load
\begin{equation}
 c_S:=b_S-\alpha\rho D_S^{1/2}\one_S.
 \label{eq:shifted-load}
\end{equation}
Whenever the restricted optimum is positive on $S$, its active coordinates
satisfy
\begin{equation}
 x[S]_S=Q_{SS}^{-1}c_S,
 \qquad
 x[S]_{V\setminus S}=0.
 \label{eq:restricted-state}
\end{equation}
The same algebra applies to an unregularized or source-shifted restricted
system after replacing $c_S$.  Thus the central object is
\begin{equation}
 G_S:=Q_{SS}^{-1}.
 \label{eq:green-operator}
\end{equation}

An iterative backend approximates $G_Sc_S$ by row operations, splittings,
coordinate updates, or a Krylov recurrence.  A response backend represents
selected actions of $G_S$ through factors, elimination messages, Schur
complements, or a hierarchy.  These are different representations of the same
operator, not different optimization problems.

\subsection{Exact bordered correction}

Let $S^+=S\cup T$ with $S\cap T=\emptyset$, and order the enlarged principal
matrix by $(S,T)$.  Put
\begin{equation}
 K_T:=Q_{TT}-Q_{TS}Q_{SS}^{-1}Q_{ST}.
 \label{eq:frontier-schur}
\end{equation}

\begin{proposition}[Exact response correction]
\label{prop:correction}
Let $x^-=x[S]$ and $x^+=x[S^+]$ solve the same fixed quadratic on their
respective faces, and put $\Delta=x^+-x^-$.  If
\begin{equation}
 h_T:=c_T-Q_{TS}x^-_S,
 \label{eq:frontier-demand}
\end{equation}
then $K_T$ is positive definite and
\begin{equation}
 \Delta_T=K_T^{-1}h_T,
 \qquad
 \Delta_S=-Q_{SS}^{-1}Q_{ST}\Delta_T.
 \label{eq:correction}
\end{equation}
Moreover, for
$\phi(x):=x^\top Qx/2-c^\top x$,
\begin{equation}
 \phi(x^-)-\phi(x^+)
 =\frac12\|\Delta\|_{Q_{S^+S^+}}^2
 =\frac12h_T^\top K_T^{-1}h_T.
 \label{eq:correction-energy}
\end{equation}
\end{proposition}

\begin{proof}
Restricted optimality gives
$(Qx^--c)_S=0$ and
$(Qx^--c)_T=-h_T$.  Subtracting the enlarged optimality equations gives
\[
 \begin{bmatrix}Q_{SS}&Q_{ST}\\Q_{TS}&Q_{TT}\end{bmatrix}
 \begin{bmatrix}\Delta_S\\\Delta_T\end{bmatrix}
 =\begin{bmatrix}0\\h_T\end{bmatrix}.
\]
Eliminating the first block proves \cref{eq:correction}.  Positive
definiteness of the principal matrix implies positive definiteness of its
Schur complement.  Completing the square around $x^+$ and substituting the
block solution proves \cref{eq:correction-energy}.
\end{proof}

The formula identifies the two costs that a literal warm start hides.  The
frontier solve uses $K_T^{-1}$, while the old-face correction uses the response
$Q_{SS}^{-1}Q_{ST}$.  Either action can be dense even when $T$ contains one
vertex.

The next result closes the numerical no-restart question whenever those
old-face responses can be applied at charged cost.  The orthogonality is not
limited to the exact correction vector: it holds between the entire lifted
Schur-frontier spaces.

\begin{theorem}[No-restart orthogonal lifted-frontier repair]
\label{thm:orthogonal-frontier-lifts}
Let $S_0\subsetneq\cdots\subsetneq S_J$ be nested faces and put
$B_j:=S_{j+1}\setminus S_j$.  Define
\begin{align}
 Z_j&:=Q_{S_jS_j}^{-1}Q_{S_jB_j},
 &K_j&:=Q_{B_jB_j}-Q_{B_jS_j}Z_j,
 \label{eq:frontier-lift-blocks}
\end{align}
and the ambient lift $\mathcal L_j:\mathbb R^{B_j}\to\mathbb R^V$ by
\begin{equation}
 (\mathcal L_jy)_{S_j}:=-Z_jy,
 \qquad
 (\mathcal L_jy)_{B_j}:=y,
 \qquad
 (\mathcal L_jy)_{V\setminus S_{j+1}}:=0.
 \label{eq:frontier-lift}
\end{equation}
Then, for arbitrary frontier vectors $u$ and $v$,
\begin{align}
 (Q\mathcal L_jv)_{S_j}&=0,
 &\|\mathcal L_jv\|_Q^2&=v^TK_jv,
 \label{eq:frontier-lift-isometry}\\
 \inner{\mathcal L_i u}{Q\mathcal L_j v}&=0
 &&(0\leq i<j<J).
 \label{eq:frontier-lift-orthogonality}
\end{align}

Let $z^j:=x[S_j]$ be the exact restricted solutions, extended by zero, and
let $y_j^*:=(z^{j+1}-z^j)_{B_j}$.  For arbitrary approximations
$\widetilde y_j\in\mathbb R^{B_j}$, set
\begin{equation}
 \widetilde z^J
 :=z^0+\sum_{j=0}^{J-1}\mathcal L_j\widetilde y_j.
 \label{eq:approximate-lifted-trajectory}
\end{equation}
Then
\begin{align}
 z^J&=z^0+\sum_{j=0}^{J-1}\mathcal L_jy_j^*,
 \label{eq:exact-lifted-trajectory}\\
 \|\widetilde z^J-z^J\|_Q^2
 &=\sum_{j=0}^{J-1}
   \|\widetilde y_j-y_j^*\|_{K_j}^2,
 \label{eq:lifted-error-pythagoras}\\
 \phi(\widetilde z^J)-\phi(z^J)
 &=\frac12\sum_{j=0}^{J-1}
   \|\widetilde y_j-y_j^*\|_{K_j}^2.
 \label{eq:lifted-objective-pythagoras}
\end{align}
If $Q$ is Stieltjes, $z^0\geq0$, and every $\widetilde y_j\geq0$, then
$\widetilde z^J\geq0$ as well.
\end{theorem}

\begin{proof}
Multiplying \cref{eq:frontier-lift} by the $S_j$ block row gives
$-Q_{S_jS_j}Z_jv+Q_{S_jB_j}v=0$.  Substitution into the block quadratic gives
\begin{equation*}
 \begin{bmatrix}-Z_jv\\v\end{bmatrix}^{\!T}
 Q_{S_{j+1}S_{j+1}}
 \begin{bmatrix}-Z_jv\\v\end{bmatrix}
 =v^TK_jv,
\end{equation*}
which proves \cref{eq:frontier-lift-isometry}.  If $i<j$, then
$\mathcal L_i u$ is supported on $S_{i+1}\subseteq S_j$, while
$(Q\mathcal L_jv)_{S_j}=0$.  This proves
\cref{eq:frontier-lift-orthogonality} for the whole pair of subspaces.

Applying \cref{prop:correction} at stage $j$ gives
$z^{j+1}-z^j=\mathcal L_jy_j^*$, proving
\cref{eq:exact-lifted-trajectory}.  Subtract the two lifted trajectories,
apply pairwise orthogonality, and then use the isometry to obtain
\cref{eq:lifted-error-pythagoras}.  Both trajectories are supported on
$S_J$, and $z^J$ is stationary on that face.  Completing the square around
$z^J$ proves \cref{eq:lifted-objective-pythagoras}.

Finally, a principal inverse of a Stieltjes matrix is nonnegative and
$Q_{S_jB_j}\leq0$, so $-Z_j\geq0$.  Thus a nonnegative frontier vector has a
nonnegative lift, proving the feasibility statement.
\end{proof}

\begin{corollary}[Charged frontier solves need no old-face restart]
\label{cor:no-restart-frontier-work}
Suppose the approximations in \cref{thm:orthogonal-frontier-lifts} satisfy
\begin{equation}
 \frac12\|\widetilde y_j-y_j^*\|_{K_j}^2\leq E_j,
 \qquad
 \sum_{j=0}^{J-1}E_j\leq E.
 \label{eq:frontier-error-budget}
\end{equation}
Then $\phi(\widetilde z^J)-\phi(z^J)\leq E$ without any numerical repair of a
previous face.  More quantitatively, suppose the $j$th preconditioned
frontier solve takes
$\widetilde O(\sqrt{\kappa_j})$ iterations; its frontier-local part costs
$O(\operatorname{cvol}(B_j))$ per iteration; and the exhaustive charge for
all response applications, old-face reads, reduced-right-hand-side
construction, boundary reporting, and lift maintenance is $U_{\rm lift}(V)$,
where $V:=\operatorname{cvol}(S_J)$.  Then all frontier repairs cost
\begin{equation}
 \widetilde O\!\left(
 U_{\rm lift}(V)+V\sqrt{\kappa_{\max}}
 \right),
 \qquad
 \kappa_{\max}:=\max_j\kappa_j.
 \label{eq:no-restart-frontier-work}
\end{equation}
\end{corollary}

\begin{proof}
The error claim is \cref{eq:lifted-objective-pythagoras}.  The batches are
disjoint, so
$\sum_j\operatorname{cvol}(B_j)\leq V$.  Sum the frontier-local iteration
charges, bound every $\sqrt{\kappa_j}$ by
$\sqrt{\kappa_{\max}}$, and add the disjoint exhaustive response charge.
Tolerance logarithms are absorbed by $\widetilde O$.
\end{proof}

This result removes repeated old-face numerical optimization from the mixed
program.  It deliberately does not make response application free:
$U_{\rm lift}$ must include every action of
$Q_{S_jS_j}^{-1}Q_{S_jB_j}$ and every sign-reporting query.  The remaining
graph-uniform problem is therefore response representation and finite-band
crossing discovery, not accumulation of frontier-solve error.

\begin{proposition}[Dual-local orthogonal correction family]
\label{prop:dual-local-orthogonal-corrections}
Let $S_0\subsetneq\cdots\subsetneq S_J$ be nested faces, let
$z^j:=x[S_j]$ be their exact restricted solutions extended by zero, and put
\begin{equation}
 B_j:=S_{j+1}\setminus S_j,
 \qquad
 \Delta^j:=z^{j+1}-z^j,
 \qquad
 q^j:=Q\Delta^j.
 \label{eq:dual-local-corrections}
\end{equation}
Then
\begin{align}
 q^j_{S_j}&=0,
 &\supp(q^j)&\subseteq B_j\cup\partial S_{j+1},
 \label{eq:dual-frontier-support}\\
 \inner{\Delta^i}{Q\Delta^j}&=0
 &&(0\leq i<j<J).
 \label{eq:nested-correction-orthogonality}
\end{align}
Consequently
\begin{equation}
 \sum_{j=0}^{J-1}\norm{\Delta^j}_Q^2
 =\norm{z^J-z^0}_Q^2
 =2\bigl(\phi(z^0)-\phi(z^J)\bigr).
 \label{eq:orthogonal-correction-telescope}
\end{equation}
If the trace is monotone, then $q^j_v\leq0$ for
$v\notin S_{j+1}$, and the normalized boundary-demand increment there is
$-q^j_v/\sqrt{d_v}\geq0$.
\end{proposition}

\begin{proof}
Both $z^j$ and $z^{j+1}$ satisfy the restricted optimality equations on
$S_j$, so subtraction gives $q^j_{S_j}=0$.  The correction is supported on
$S_{j+1}$; sparsity of $Q$, followed by the old-face cancellation, proves
the support statement.  If $i<j$, then $\Delta^i$ is supported on
$S_{i+1}\subseteq S_j$, where $q^j$ vanishes.  This proves
\cref{eq:nested-correction-orthogonality}.  The increments sum to
$z^J-z^0$, so Pythagoras in the $Q$ inner product proves the first equality
in \cref{eq:orthogonal-correction-telescope}; summing
\cref{eq:correction-energy} proves the second.  Under monotonicity,
$\Delta^j\geq0$, and every row outside $S_{j+1}$ contains only nonpositive
off-diagonal coefficients against its support.  Hence $q_v^j\leq0$ and the
last claim follows by subtracting the two boundary demands.
\end{proof}

This proposition answers the basic orthogonality--locality question at the
algebraic level.  Orthogonality survives nested support growth, and the dual
signature $q^j$ is frontier-local even when the primal vector $\Delta^j$ is
dense.  The unresolved algorithmic step is to produce or query that signature
from the implicit response without first materializing the dense primal
correction.

\begin{corollary}[Positive dual response transfer]
\label{cor:positive-dual-transfer}
In \cref{prop:correction}, put $R:=V\setminus(S\cup T)$ and define the
post-elimination cross block
\begin{equation}
 C_{RT}^{(S)}
 :=Q_{RT}-Q_{RS}Q_{SS}^{-1}Q_{ST}.
 \label{eq:post-elimination-cross-block}
\end{equation}
Then the increase in exterior demand caused by the exact expansion is
\begin{equation}
 y_R
 :=-Q_{R,S\cup T}\Delta_{S\cup T}
 =-C_{RT}^{(S)}K_T^{-1}h_T.
 \label{eq:positive-dual-transfer}
\end{equation}
The Schur complement obtained by eliminating $S$ is Stieltjes, so
$C_{RT}^{(S)}\leq0$ and $K_T^{-1}\geq0$.  Hence the transfer operator
\begin{equation}
 \mathcal R_{R\leftarrow T}^{(S)}
 :=-C_{RT}^{(S)}K_T^{-1}
 \label{eq:dual-response-operator}
\end{equation}
is entrywise nonnegative.  In particular, $h_T\geq0$ implies $y_R\geq0$.
\end{corollary}

\begin{proof}
Substitute
$\Delta_S=-Q_{SS}^{-1}Q_{ST}\Delta_T$ and
$\Delta_T=K_T^{-1}h_T$ from \cref{eq:correction} into
$-Q_{RS}\Delta_S-Q_{RT}\Delta_T$.  Principal Schur complements of a
positive-definite Stieltjes matrix are positive-definite Stieltjes matrices;
their off-diagonal blocks are nonpositive and their inverses are
nonnegative.
\end{proof}

Thus a nested exact-face algorithm can be written entirely as positive dual
transport: the current violating demand $h_T$ is sent through
$\mathcal R_{R\leftarrow T}^{(S)}$, the resulting threshold crossings form
the next batch, and the dense primal correction is postponed.  The
graph-uniform finite-band problem is precisely to apply this positive transfer
implicitly and report all crossings at output-sensitive cost.

\begin{corollary}[Associative batching of dual transfers]
\label{cor:associative-dual-batching}
Under the hypotheses and notation of
\cref{prop:dual-local-orthogonal-corrections}, fix
$0\leq a<b\leq J$ and put
\begin{equation*}
 \Delta^{[a,b)}:=z^b-z^a,
 \qquad
 q^{[a,b)}:=Q\Delta^{[a,b)}.
\end{equation*}
Then
\begin{align}
 q^{[a,b)}
 &=\sum_{j=a}^{b-1}q^j,
 &q^{[a,b)}_{S_a}&=0,
 &\supp(q^{[a,b)})
 &\subseteq(S_b\setminus S_a)\cup\partial S_b.
 \label{eq:batched-dual-signature}
\end{align}
For $R:=V\setminus S_b$, the cumulative exterior-demand increase is
\begin{equation}
 -q_R^{[a,b)}
 =\sum_{j=a}^{b-1}(-q_R^j)\geq0.
 \label{eq:batched-exterior-demand}
\end{equation}
Moreover, writing $Q/U$ for the Schur complement after eliminating $U$,
\begin{equation}
 (Q/S_a)/(S_b\setminus S_a)=Q/S_b.
 \label{eq:schur-quotient-associativity}
\end{equation}
Thus replacing any consecutive interval of exact light events by its single
aggregate transfer preserves the endpoint solution, every exterior demand,
and the Schur operator seen by all future events.
\end{corollary}

\begin{proof}
The primal corrections telescope, and multiplication by $Q$ gives the first
identity.  Both endpoint solutions satisfy the restricted equations on
$S_a$, so the same old-face cancellation used in
\cref{prop:dual-local-orthogonal-corrections} proves the zero and support
statements.  Restricting the telescoping identity to $R$ proves
\cref{eq:batched-exterior-demand}; monotonicity gives its sign.  Finally,
block Gaussian elimination is associative: eliminating the variables in
$S_a$ and then those in $S_b\setminus S_a$ produces the same block factorization
and residual quadratic as eliminating $S_b$ at once.  This proves
\cref{eq:schur-quotient-associativity} and the endpoint claims.
\end{proof}

The corollary removes a possible logical obstruction to a heavy/light
hierarchy: delayed light transfers may be merged exactly, rather than
approximating a different active-set trajectory.  What remains is the data
structural question of representing the aggregate positive transfer and
finding its finite-band crossings without replaying all constituent events.

\section{A response-preconditioned state}
\label{sec:state}

At an epoch, split the current face as
\begin{equation}
 S=A\mathbin{\dot\cup}F,
 \label{eq:anchor-frontier}
\end{equation}
where $A$ is a settled anchor and $F$ is a frontier buffer.  The backend stores:

\begin{enumerate}[label=(\roman*)]
\item a persistent response representation for selected actions of
      $Q_{AA}^{-1}$;
\item the exposed blocks $Q_{AF}$ and $Q_{FF}$;
\item a frontier Schur operator
      $Q_{FF}-Q_{FA}Q_{AA}^{-1}Q_{AF}$, exact or certified approximate;
\item an implicit numerical state and error certificate on $S$;
\item a boundary dictionary whose keys are intervals, not uncertified point
      estimates;
\item a work ledger charging adjacency exposure, old-face reads, response
      updates, rekeys, iterative row operations, and final writes.
\end{enumerate}

The controller requests four conceptual operations.

\begin{description}[leftmargin=2.8cm,style=nextline]
\item[\textsc{Expand}$(T)$]
Expose and add a support-safe batch, updating the anchor--frontier
representation without materializing every old coordinate.
\item[\textsc{BoundaryBounds}]
Return all certified violations and certify that every unreported boundary
coordinate lies below the admission threshold or inside the declared
transition band.
\item[\textsc{Repair}$(\delta)$]
Reduce numerical error on the current face using an exact solve or a
preconditioned iterative method.
\item[\textsc{Finalize}]
Materialize the terminal vector once and return its note-scoped correctness
certificate.
\end{description}

The completeness clause in \textsc{BoundaryBounds} is essential.  A data
structure that answers the value at one requested vertex is not yet a local
active-set backend because it may not find an unqueried violating vertex.

\section{Certified finite-band boundary intervals}
\label{sec:intervals}

Fix a current face $S$ and a boundary vertex $v$.  Let
\begin{equation}
 a_v:=Q_{Sv},
 \qquad
 \chi_v^2:=a_v^\top Q_{SS}^{-1}a_v.
 \label{eq:boundary-leverage}
\end{equation}
The scalar $\chi_v$ is a response leverage.  It is closely related to the
energy needed to transmit a unit boundary demand through the active face.

\begin{lemma}[Energy-to-boundary interval]
\label{lem:boundary-interval}
Suppose $x_S$ is the exact restricted solution and $\widetilde x_S$ satisfies
\begin{equation}
 \|\widetilde x_S-x_S\|_{Q_{SS}}\leq\delta.
 \label{eq:energy-error}
\end{equation}
For every queried boundary vertex $v$,
\begin{equation}
 \left|
 Q_{vS}(\widetilde x_S-x_S)
 \right|
 \leq \chi_v\delta.
 \label{eq:boundary-error}
\end{equation}
Consequently a point estimate of any boundary demand affine in $Q_{vS}x_S$
becomes a certified interval after adding and subtracting $\chi_v\delta$.
\end{lemma}

\begin{proof}
Write the left side as
\[
 \left|
 (Q_{SS}^{-1/2}a_v)^\top
 Q_{SS}^{1/2}(\widetilde x_S-x_S)
 \right|.
\]
Cauchy--Schwarz and \cref{eq:boundary-leverage,eq:energy-error} prove the
claim.
\end{proof}

\begin{corollary}[Graph-universal leverage bound]
\label{cor:boundary-leverage-one}
Every response leverage in \cref{eq:boundary-leverage} satisfies
\begin{equation}
 \chi_v^2\leq Q_{vv}-\alpha=\frac{1-\alpha}{2}<1.
 \label{eq:boundary-leverage-one}
\end{equation}
Consequently
\begin{equation}
 \frac{\left|Q_{vS}(\widetilde x_S-x_S)\right|}{\sqrt{d_v}}
 \leq\delta
 \label{eq:normalized-boundary-error}
\end{equation}
under \cref{eq:energy-error}, without estimating a coordinate-specific
leverage.
\end{corollary}

\begin{proof}
The scalar Schur complement is
$K_v:=Q_{vv}-Q_{vS}Q_{SS}^{-1}Q_{Sv}$.  Its variational characterization and
$Q\succeq\alpha I$ give
\begin{equation*}
 K_v
 =\min_y\begin{bmatrix}y\\1\end{bmatrix}^{\!T}
 Q_{S\cup\{v\},S\cup\{v\}}
 \begin{bmatrix}y\\1\end{bmatrix}
 \geq\alpha.
\end{equation*}
The shared normalized matrix has $Q_{vv}=(1+\alpha)/2$, proving the first
claim.  Apply \cref{lem:boundary-interval} and $d_v\geq1$ for the displayed
normalized error bound.
\end{proof}

If the lower endpoint exceeds the gate, admission is safe.  If the upper
endpoint is below the gate, rejection is safe.  Only intervals intersecting
the transition band require refinement.  The lemma proves correctness of one
queried interval; it does not bound the work needed to find all ambiguous or
violating vertices.  That reporting problem is the main open response theorem.

The graph-universal leverage is intentionally oblivious to where a correction
was generated.  A frontier correction has much more structure: its dual
source is supported on the newly admitted batch.  The next lemma converts that
support into a trace bound on all boundary sensitivities.

\begin{lemma}[Source-aware leverage packing]
\label{lem:source-aware-leverage-packing}
Fix a face $S$, an exterior set $W\subseteq V\setminus S$, and a nonempty
source set $B\subseteq S$.  Let $E_B$ be the coordinate embedding from
$\mathbb R^B$ into $\mathbb R^S$ and put
\begin{equation}
 H_B:=E_B^TQ_{SS}^{-1}E_B.
 \label{eq:source-green-block}
\end{equation}
For $v\in W$, define the source-aware leverage
\begin{equation}
 (\chi_v^{B})^2
 :=Q_{vS}Q_{SS}^{-1}E_BH_B^{-1}
       E_B^TQ_{SS}^{-1}Q_{Sv}.
 \label{eq:source-aware-leverage}
\end{equation}
Then every correction of the form
\begin{equation}
 \Delta=Q_{SS}^{-1}E_Bh
 \label{eq:source-supported-correction}
\end{equation}
satisfies
\begin{equation}
 |Q_{vS}\Delta|
 \leq\chi_v^B\|\Delta\|_{Q_{SS}},
 \qquad v\in W,
 \label{eq:source-aware-interval}
\end{equation}
and the complete exterior leverage ledger obeys
\begin{equation}
 \sum_{v\in W}(\chi_v^B)^2\leq |B|.
 \label{eq:source-aware-leverage-sum}
\end{equation}
Consequently, for arbitrary positive normalized interval margins $m_v$, the
set
\begin{equation*}
 \mathcal A
 :=\left\{v\in W:
   \frac{\chi_v^B\|\Delta\|_{Q_{SS}}}{\sqrt{d_v}}
   \geq m_v
 \right\}
\end{equation*}
has the a priori packing bound
\begin{equation}
 \sum_{v\in\mathcal A}d_vm_v^2
 \leq |B|\,\|\Delta\|_{Q_{SS}}^2.
 \label{eq:source-aware-ambiguous-packing}
\end{equation}
For a common margin $m$, this gives
$\vol(\mathcal A)\leq |B|\|\Delta\|_{Q_{SS}}^2/m^2$.
Every nested exact correction in
\cref{prop:dual-local-orthogonal-corrections} has
\cref{eq:source-supported-correction} on $S=S_{j+1}$ with $B=B_j$.
\end{lemma}

\begin{proof}
The Green block $H_B$ is positive definite.  Define
\begin{equation*}
 P_B:=Q_{SS}^{-1/2}E_BH_B^{-1}E_B^TQ_{SS}^{-1/2}.
\end{equation*}
It is the orthogonal projector onto
$\operatorname{range}(Q_{SS}^{-1/2}E_B)$ and has trace $|B|$.
Also put $R:=Q_{WS}Q_{SS}^{-1/2}$.  Schur complementation of the principal
matrix on $S\cup W$, followed by $Q\preceq I$, gives
\begin{equation}
 RR^T=Q_{WS}Q_{SS}^{-1}Q_{SW}\preceq Q_{WW}\preceq I.
 \label{eq:exterior-response-contraction}
\end{equation}
Thus $R^TR\preceq I$ as well.  The definition
\cref{eq:source-aware-leverage} is the $v$th diagonal of $RP_BR^T$, so
\begin{equation*}
 \sum_{v\in W}(\chi_v^B)^2
 =\operatorname{tr}(P_BR^TR)
 \leq\operatorname{tr}(P_B)=|B|,
\end{equation*}
proving \cref{eq:source-aware-leverage-sum}.

For \cref{eq:source-supported-correction}, let
$w_v:=E_B^TQ_{SS}^{-1}Q_{Sv}$.  Then
\begin{equation*}
 Q_{vS}\Delta=w_v^Th,
 \qquad
 \|\Delta\|_{Q_{SS}}^2=h^TH_Bh.
\end{equation*}
Cauchy--Schwarz in the $H_B$ inner product gives
$|w_v^Th|^2\leq(w_v^TH_B^{-1}w_v)(h^TH_Bh)$, which is exactly
\cref{eq:source-aware-interval}.  Every $v\in\mathcal A$ satisfies
$d_vm_v^2\leq(\chi_v^B)^2\|\Delta\|_{Q_{SS}}^2$; summing and applying
\cref{eq:source-aware-leverage-sum} proves
\cref{eq:source-aware-ambiguous-packing}.  Finally,
\cref{eq:dual-frontier-support} says that
$Q_{S_{j+1}S_{j+1}}\Delta^j$ is supported on $B_j$, which gives the last
claim by solving the enlarged principal system.
\end{proof}

The source-aware ledger is an a priori statement: it bounds the volume whose
certified interval can be wide at a declared margin before the actual exterior
change is computed.  The following identity shows that its masses are not new
matrix objects.  They are exactly the diagonal loss of the terminal Schur
operator across the corresponding elimination.

\begin{theorem}[Source-aware leverage equals Schur-diagonal loss]
\label{thm:source-leverage-diagonal-loss}
Let $A\subsetneq T=A\mathbin{\dot\cup}B$ and
$W\subseteq V\setminus T$.  Write
\begin{equation*}
 L^A:=Q/A,
 \qquad
 L^T:=Q/T
\end{equation*}
for the terminal Schur complements, restricted to their surviving
coordinates.  On the enlarged face $T$, let $E_B$ embed $\mathbb R^B$ into
$\mathbb R^T$, put $H_B:=E_B^TQ_{TT}^{-1}E_B$, and define
\begin{equation}
 (\chi_v^{B\mid A})^2
 :=Q_{vT}Q_{TT}^{-1}E_BH_B^{-1}
      E_B^TQ_{TT}^{-1}Q_{Tv},
 \qquad v\in W.
 \label{eq:conditional-source-leverage}
\end{equation}
Then
\begin{equation}
 (\chi_v^{B\mid A})^2
 =L^A_{vB}(L^A_{BB})^{-1}L^A_{Bv}
 =L^A_{vv}-L^T_{vv}.
 \label{eq:source-leverage-diagonal-loss}
\end{equation}
Consequently every group $C\subseteq W$ has the exact additive mass
\begin{equation}
 Z_{B\mid A}(C)
 :=\sum_{v\in C}\frac{(\chi_v^{B\mid A})^2}{d_v}
 =\operatorname{tr}\!\left[
  D_C^{-1}\bigl(L^A_{CC}-L^T_{CC}\bigr)
 \right].
 \label{eq:group-diagonal-loss}
\end{equation}

More generally, let
$S_0\subsetneq\cdots\subsetneq S_J$ and
$B_j:=S_{j+1}\setminus S_j$.  If $v\notin S_k$, then its source-aware
sensitivity has the temporal telescope
\begin{equation}
 \sum_{j=0}^{k-1}(\chi_v^{B_j\mid S_j})^2
 =(Q/S_0)_{vv}-(Q/S_k)_{vv}
 \leq Q_{vv}-\alpha
 =\frac{1-\alpha}{2}.
 \label{eq:source-leverage-temporal-telescope}
\end{equation}
For every fixed $C\subseteq V\setminus S_k$, the normalized group masses
telescope as well:
\begin{equation}
 \sum_{j=0}^{k-1}Z_{B_j\mid S_j}(C)
 =\operatorname{tr}\!\left[
 D_C^{-1}\bigl((Q/S_0)_{CC}-(Q/S_k)_{CC}\bigr)
 \right].
 \label{eq:group-diagonal-temporal-telescope}
\end{equation}
Thus every still-exterior label has a graph-universal total squared
source-aware sensitivity budget, independent of the number of preceding
expansions.
\end{theorem}

\begin{proof}
Order $T$ as $(A,B)$ and put $K_B:=L^A_{BB}$.  The block inverse formula gives
\begin{equation*}
 H_B=K_B^{-1},
 \qquad
 Q_{vT}Q_{TT}^{-1}E_B=L^A_{vB}K_B^{-1}.
\end{equation*}
Substitution in \cref{eq:conditional-source-leverage} proves the first
equality in \cref{eq:source-leverage-diagonal-loss}.  Schur associativity
gives
\begin{equation*}
 L^T_{WW}
 =L^A_{WW}-L^A_{WB}K_B^{-1}L^A_{BW},
\end{equation*}
whose $v$th diagonal proves the second equality.  Summing these diagonal
identities with weights $1/d_v$ proves \cref{eq:group-diagonal-loss}.

Apply the same identity at every consecutive pair $(S_j,S_{j+1})$ and
telescope.  Since $Q\succeq\alpha I$, the variational characterization of a
Schur complement gives $Q/S_k\succeq\alpha I$.  Also
$(Q/S_0)_{vv}\leq Q_{vv}$.  The shared PageRank matrix has
$Q_{vv}=(1+\alpha)/2$, proving
\cref{eq:source-leverage-temporal-telescope}.  Summing the same termwise
diagonal telescope over $v\in C$ with weights $1/d_v$ proves
\cref{eq:group-diagonal-temporal-telescope}.
\end{proof}

The identity replaces source-conditioned response matrices by one monotone
scalar field: terminal Schur diagonals.  It does not make those diagonals
free, but a dynamic implementation may now obtain every group mass by a
before--after aggregate diagonal query.  In particular, no multiplicative
accuracy is needed when the true decrement is far below the current search
scale.  The next lemma first states the clean multiplicative hierarchy, and
the following corollary weakens its oracle to a one-sided additive guarantee.

\begin{lemma}[Rank-sensitive hierarchical locator]
\label{lem:source-aware-hierarchical-locator}
Under the hypotheses of
\cref{lem:source-aware-leverage-packing}, let $\mathcal H$ be a rooted
partition hierarchy whose leaves are the vertices of $W$, whose arity is at
most $b$, and whose depth is $L$.  For every node $C\subseteq W$, define its
degree-normalized source-aware mass by
\begin{equation}
 Z_B(C):=\sum_{v\in C}\frac{(\chi_v^B)^2}{d_v}.
 \label{eq:source-aware-group-mass}
\end{equation}
Suppose a group query returns a certified multiplicative upper estimate
\begin{equation}
 Z_B(C)\leq\widehat Z_B(C)\leq\gamma Z_B(C)
 \label{eq:source-aware-group-oracle}
\end{equation}
for some $\gamma\geq1$.  Fix a correction
$\Delta=Q_{SS}^{-1}E_Bh$ with energy radius
$E:=\|\Delta\|_{Q_{SS}}>0$ and a common normalized margin $m>0$.

Starting at the root, discard a queried node $C$ when
$\widehat Z_B(C)<m^2/E^2$ and otherwise query its children.  At every reached
leaf, evaluate the exact leaf predicate.  This search returns precisely the
potentially ambiguous set
\begin{equation*}
 \mathcal A_m
 :=\left\{v\in W:
     \frac{\chi_v^B E}{\sqrt{d_v}}\geq m
   \right\},
\end{equation*}
uses at most
\begin{equation}
 1+bL\left(1+\gamma|B|\frac{E^2}{m^2}\right)
 \label{eq:source-aware-group-query-count}
\end{equation}
group queries, and has output charge
\begin{equation}
 \vol(\mathcal A_m)\leq |B|\frac{E^2}{m^2}.
 \label{eq:source-aware-locator-output}
\end{equation}
Thus, if group queries and exact reached-leaf tests cost polylogarithmic work,
the search is sensitive to the packed ambiguity budget rather than to the
ambient boundary size, up to the hierarchy depth and approximation factor.
\end{lemma}

\begin{proof}
The masses in \cref{eq:source-aware-group-mass} are nonnegative and additive
over the children of every hierarchy node.  Since $d_v\geq1$,
\cref{eq:source-aware-leverage-sum} gives
\begin{equation}
 Z_B(W)\leq\sum_{v\in W}(\chi_v^B)^2\leq |B|.
 \label{eq:source-aware-normalized-mass-total}
\end{equation}
If $v\in\mathcal A_m$, then every ancestor $C$ of $v$ satisfies
$Z_B(C)\geq Z_B(\{v\})\geq m^2/E^2$, and therefore cannot be discarded by
the upper estimate in \cref{eq:source-aware-group-oracle}.  Conversely, the
exact predicate at a reached leaf removes every false candidate, proving
correctness.

At any fixed depth the retained nodes are disjoint.  A retained node satisfies
$\widehat Z_B(C)\geq m^2/E^2$, and the upper side of
\cref{eq:source-aware-group-oracle} implies
$Z_B(C)\geq m^2/(\gamma E^2)$.  Hence
\cref{eq:source-aware-normalized-mass-total} bounds the number of retained
nodes at that depth by $\gamma|B|E^2/m^2$.  Querying at most $b$ children per
retained node over $L$ levels, together with the root and harmless integer
rounding, proves \cref{eq:source-aware-group-query-count}.  The output-volume
claim is the common-margin case of
\cref{eq:source-aware-ambiguous-packing}.
\end{proof}

\begin{corollary}[Margin-adaptive additive group oracle]
\label{cor:additive-group-locator}
In \cref{lem:source-aware-hierarchical-locator}, put
$\tau:=m^2/E^2$ and replace the multiplicative oracle by certified estimates
satisfying
\begin{equation}
 Z_B(C)\leq\widehat Z_B(C)
 \leq\gamma Z_B(C)+\beta\tau,
 \qquad 0\leq\beta<1.
 \label{eq:additive-group-oracle}
\end{equation}
Use the same rule: discard $C$ when $\widehat Z_B(C)<\tau$, and otherwise
descend.  The search still returns exactly $\mathcal A_m$, and it uses at
most
\begin{equation}
 1+bL\left(
  1+\frac{\gamma}{1-\beta}|B|\frac{E^2}{m^2}
 \right)
 \label{eq:additive-group-query-count}
\end{equation}
group queries.  The oracle may therefore have additive error proportional to
the current pruning threshold; it need not resolve arbitrarily small group
masses multiplicatively.  By
\cref{thm:source-leverage-diagonal-loss}, the same statement holds when each
$Z_B(C)$ is supplied as the weighted terminal-Schur diagonal decrement in
\cref{eq:group-diagonal-loss}.
\end{corollary}

\begin{proof}
If $v\in\mathcal A_m$, every ancestor has
$Z_B(C)\geq Z_B(\{v\})\geq\tau$, so the certified upper estimate cannot
discard it.  Exact leaf tests again remove all false candidates.  Conversely,
every retained node satisfies
\begin{equation*}
 \tau\leq\widehat Z_B(C)
 \leq\gamma Z_B(C)+\beta\tau,
\end{equation*}
and hence $Z_B(C)\geq(1-\beta)\tau/\gamma$.  Retained nodes at one depth are
disjoint, while $Z_B(W)\leq|B|$.  The level count and the child-query charge
are therefore the same as in
\cref{lem:source-aware-hierarchical-locator}, with the factor
$\gamma/(1-\beta)$, proving
\cref{eq:additive-group-query-count}.
\end{proof}

The group mass has the exact trace form
\begin{equation*}
 Z_B(C)
 =\operatorname{tr}\!\left(
 P_BR_C^TD_C^{-1}R_C
 \right),
 \qquad
 R_C:=Q_{CS}Q_{SS}^{-1/2},
\end{equation*}
where $D_C:=\operatorname{diag}(d_v:v\in C)$.  Thus the remaining P0
primitive is no longer an unspecified boundary scan: it is a dynamic,
constant-factor group-trace oracle for these source-conditioned response
rows, together with exact evaluation only at the reached leaves.  The lemma
does not construct that oracle on a general nonequitable cyclic core.

The next reduction shows that the trace oracle does not require one response
query per hierarchy node.  A logarithmic batch of source-normalized random
responses estimates every fixed node simultaneously, including nodes chosen
adaptively by the subsequent branch-and-bound search.

\begin{lemma}[Group traces from logarithmically many response sketches]
\label{lem:group-trace-response-sketch}
Under the hypotheses of
\cref{lem:source-aware-hierarchical-locator}, let $\mathcal C$ be the set of
$M$ nodes of the fixed hierarchy.  Fix $\xi,\delta\in(0,1)$, draw
$r$ independent standard Gaussian vectors
$g_1,\ldots,g_r\in\mathbb R^B$, and define the exterior response sketches
\begin{equation}
 s^{(t)}
 :=D_W^{-1/2}Q_{WS}Q_{SS}^{-1}E_BH_B^{-1/2}g_t,
 \qquad 1\leq t\leq r.
 \label{eq:source-normalized-response-sketch}
\end{equation}
For every hierarchy node $C$, put
\begin{equation}
 \overline Z_B(C)
 :=\frac1r\sum_{t=1}^r\|s_C^{(t)}\|_2^2,
 \qquad
 \widehat Z_B(C):=\frac{\overline Z_B(C)}{1-\xi}.
 \label{eq:sketched-group-trace}
\end{equation}
There is a universal constant $c_{\rm sk}>0$ such that, if
\begin{equation}
 r\geq c_{\rm sk}\xi^{-2}\log(2M/\delta),
 \label{eq:group-trace-sketch-count}
\end{equation}
then, with probability at least $1-\delta$, simultaneously for every
$C\in\mathcal C$,
\begin{equation}
 Z_B(C)\leq\widehat Z_B(C)
 \leq\frac{1+\xi}{1-\xi}Z_B(C).
 \label{eq:sketched-group-oracle}
\end{equation}
In particular, $\xi=1/3$ supplies the oracle in
\cref{eq:source-aware-group-oracle} with $\gamma=2$ using
$r=O(\log(M/\delta))$ response sketches.
\end{lemma}

\begin{proof}
For a node $C$, define
\begin{equation*}
 T_C:=D_C^{-1/2}Q_{CS}Q_{SS}^{-1}E_BH_B^{-1/2}.
\end{equation*}
Cyclicity of trace and the definition of $P_B$ give
\begin{equation*}
 \|T_C\|_F^2
 =\operatorname{tr}\!\left(P_BR_C^TD_C^{-1}R_C\right)
 =Z_B(C).
\end{equation*}
Moreover, $s_C^{(t)}=T_Cg_t$, so
$\overline Z_B(C)=r^{-1}\sum_tg_t^T(T_C^TT_C)g_t$.  The standard Gaussian
quadratic-form moment-generating-function bound gives, for every positive
semidefinite matrix $A$ and $0<\xi<1$,
\begin{equation*}
 \Pr\!\left[
  \left|\frac1r\sum_{t=1}^rg_t^TAg_t-\operatorname{tr}(A)\right|
  >\xi\operatorname{tr}(A)
 \right]
 \leq2\exp(-c r\xi^2)
\end{equation*}
for a universal $c>0$.  Apply this with $A=T_C^TT_C$ and take a union bound
over all $M$ hierarchy nodes.  On the resulting event,
$(1-\xi)Z_B(C)\leq\overline Z_B(C)\leq
(1+\xi)Z_B(C)$ simultaneously.  Rescaling the estimate proves
\cref{eq:sketched-group-oracle}.  Nodes of zero mass have $T_C=0$ and obey
the same statement deterministically.
\end{proof}

Once the leaf values of the sketches are available, all quantities in
\cref{eq:sketched-group-trace} are additive and can be aggregated bottom-up.
The remaining challenge is therefore narrower still: maintain a polylogarithmic
batch of the signed response vectors in
\cref{eq:source-normalized-response-sketch}, and their squared subtree sums,
under terminal additions without materializing every exterior coordinate.
The lemma is a static randomized reduction and does not claim that dynamic
maintenance has local work.

Orthogonal frontier lifts make the sketches causal and mergeable across an
adaptive support trace.  The important choice is to accumulate the
nonnegative per-batch squared masses, rather than squaring one signed sketch
that has been reused to choose later batches.  Fresh randomness after each
batch avoids an adaptive-sketching gap.

\begin{theorem}[Causal append-only orthogonal response sketches]
\label{thm:causal-orthogonal-response-sketches}
Use the nested faces, Schur complements, and lifts of
\cref{thm:orthogonal-frontier-lifts}, and define the normalized lifted basis
\begin{equation}
 U_j:=\mathcal L_jK_j^{-1/2}\in\mathbb R^{V\times B_j}.
 \label{eq:normalized-frontier-basis}
\end{equation}
The bordered inverse identity also gives
\begin{equation}
 E_{B_j}^TQ_{S_{j+1}S_{j+1}}^{-1}E_{B_j}=K_j^{-1},
 \qquad
 U_j
 =Q_{S_{j+1}S_{j+1}}^{-1}E_{B_j}K_j^{1/2},
 \label{eq:normalized-frontier-source-identity}
\end{equation}
where the second identity is extended by zero outside $S_{j+1}$.  Thus its
exterior row norms are exactly the source-aware leverages of
\cref{eq:source-aware-leverage}.
Then
\begin{equation}
 U_i^TQU_j=
 \begin{cases}
 I_{B_j},&i=j,\\
 0,&i\neq j.
 \end{cases}
 \label{eq:orthonormal-frontier-bases}
\end{equation}
For an exterior vertex $v\notin S_{k+1}$, put
\begin{equation}
 \lambda_{jv}^2
 :=\frac{\|e_v^TQU_j\|_2^2}{d_v},
 \qquad
 \Lambda_{kv}^2:=\sum_{j=0}^k\lambda_{jv}^2.
 \label{eq:cumulative-frontier-leverage}
\end{equation}
If $Q\preceq I$ and $R_k:=\sum_{j=0}^k|B_j|$, then
\begin{equation}
 \sum_{v\notin S_{k+1}}\Lambda_{kv}^2\leq R_k.
 \label{eq:cumulative-frontier-leverage-packing}
\end{equation}
Moreover, every accumulated lifted error
$e=\sum_{j=0}^kU_ja_j$ satisfies
\begin{equation}
 \frac{|(Qe)_v|}{\sqrt{d_v}}
 \leq\Lambda_{kv}\|e\|_Q.
 \label{eq:cumulative-frontier-interval}
\end{equation}

Now allow the exact support batches and current boundary hierarchy to depend
on every earlier random choice.  After $B_j$, $K_j$, and the current hierarchy
are fixed, draw fresh independent standard Gaussian vectors
$g_{j,1},\ldots,g_{j,r_j}\in\mathbb R^{B_j}$ and define the per-batch group
increment
\begin{equation}
 \overline Z_j(C)
 :=\frac1{r_j}\sum_{t=1}^{r_j}
 \left\|D_C^{-1/2}Q_{CV}U_jg_{j,t}\right\|_2^2.
 \label{eq:append-only-group-mass-increment}
\end{equation}
Let $M_j$ bound the number of live hierarchy nodes and leaves at event $j$,
fix $\xi,\delta\in(0,1)$, and set
\begin{equation}
 \delta_j:=\frac{6\delta}{\pi^2(j+1)^2}.
 \label{eq:adaptive-sketch-failure-budget}
\end{equation}
For
\begin{equation}
 r_j\geq c_{\rm sk}\xi^{-2}\log(2M_j/\delta_j),
 \label{eq:adaptive-sketch-count}
\end{equation}
with probability at least $1-\delta$, simultaneously at every event, every
stored leaf and every current group has cumulative estimate
\begin{equation}
 \sum_{i\leq j}Z_i(C)
 \leq
 \frac1{1-\xi}\sum_{i\leq j}\overline Z_i(C)
 \leq
 \frac{1+\xi}{1-\xi}\sum_{i\leq j}Z_i(C),
 \label{eq:adaptive-cumulative-group-oracle}
\end{equation}
where
$Z_i(C):=\|D_C^{-1/2}Q_{CV}U_i\|_F^2$ and contributions are retained only
while their leaf remains exterior.  If a vertex $v$ first enters the boundary
after batch $B_j$, then $\lambda_{iv}=0$ for every $i<j$.  Its accumulator
therefore starts at zero and no earlier response sketch is replayed.  Put
$N:=1+\operatorname{cvol}(S_J)$.  Since every batch is nonempty,
$J\leq |S_J|$; if $M_j\leq N^{O(1)}$, the total number of sampled response
right-hand sides is $O_\xi(J\log(NJ/\delta))$.
\end{theorem}

\begin{proof}
The bottom-right block of the bordered inverse is $K_j^{-1}$, proving the
first identity in \cref{eq:normalized-frontier-source-identity}.  Also,
$Q_{S_{j+1}S_{j+1}}\mathcal L_jy=E_{B_j}K_jy$; substitute
$y=K_j^{-1/2}a$ and solve the principal system to obtain the second identity.
The isometry and cross-space orthogonality in
\cref{eq:frontier-lift-isometry,eq:frontier-lift-orthogonality} give
\cref{eq:orthonormal-frontier-bases}.  Let
$U_{\leq k}:=[U_0\ \cdots\ U_k]$.  Then
$Q^{1/2}U_{\leq k}$ has $R_k$ orthonormal columns.  For
$W:=V\setminus S_{k+1}$,
\begin{align*}
 \sum_{v\in W}\Lambda_{kv}^2
 &=\|D_W^{-1/2}Q_{WV}U_{\leq k}\|_F^2\\
 &\leq\|Q_{WV}U_{\leq k}\|_F^2
 \leq R_k,
\end{align*}
because $d_v\geq1$ and
$\|E_W^TQ^{1/2}\|_2\leq\|Q^{1/2}\|_2\leq1$.  Writing
$a:=(a_0,\ldots,a_k)$, orthonormality gives
$\|e\|_Q=\|a\|_2$; rowwise Cauchy--Schwarz proves
\cref{eq:cumulative-frontier-interval}.

At event $j$, condition on the complete history through the choice of
$U_j$ and the current hierarchy.  They are fixed before the fresh Gaussian
vectors are drawn.  Apply the quadratic-form concentration argument from
\cref{lem:group-trace-response-sketch} with failure probability $\delta_j$
to all current nodes and leaves.  Thus, conditionally on every possible
history, the event-$j$ estimates lie between
$(1-\xi)Z_j(C)$ and $(1+\xi)Z_j(C)$ except with probability $\delta_j$.
Since $\sum_{j\geq0}\delta_j=\delta$, a conditional union bound proves the
simultaneous claim despite adaptive future batches.  Summing nonnegative
per-batch estimates preserves both multiplicative inequalities; group
estimates are sums of their retained leaf estimates.

Finally, suppose $v$ first becomes adjacent to the active face after $B_j$.
For $i<j$, the support of $U_i$ lies in
$S_{i+1}\subseteq S_j$, while $v$ has no edge to $S_j$.  Hence
$e_v^TQU_i=0$, proving causal zero initialization.
Finally, $J\leq|S_J|$ follows from disjoint nonempty batches.  Under the
polynomial hierarchy-size assumption,
\cref{eq:adaptive-sketch-count,eq:adaptive-sketch-failure-budget} gives
$r_j=O_\xi(\log(NJ/\delta))$ uniformly; summing over the events proves the
sample count.
\end{proof}

\begin{corollary}[Cumulative sketches collapse to diagonal loss]
\label{cor:cumulative-sketch-diagonal-loss}
In \cref{thm:causal-orthogonal-response-sketches}, fix $k$ and a hierarchy
node $C\subseteq V\setminus S_{k+1}$.  The exact cumulative mass being
estimated has the deterministic representation
\begin{equation}
 \sum_{j=0}^k Z_j(C)
 =\operatorname{tr}\!\left[
 D_C^{-1}\bigl((Q/S_0)_{CC}-(Q/S_{k+1})_{CC}\bigr)
 \right].
 \label{eq:cumulative-sketch-diagonal-loss}
\end{equation}
Therefore a dynamic oracle that directly supplies the one-sided additive
estimates in \cref{eq:additive-group-oracle} for these diagonal losses can
replace all per-batch Gaussian response sketches.  Combined with
\cref{cor:additive-group-locator}, it gives the same no-false-negative
branch-and-bound search with output-sensitive query count.  No source vector,
old sketch, or past batch must be replayed.
\end{corollary}

\begin{proof}
The normalized lifted basis $U_j$ is the source-normalized response basis for
$B_j$ by \cref{eq:normalized-frontier-source-identity}.  Hence
\begin{equation*}
 Z_j(C)
 =\sum_{v\in C}\frac{(\chi_v^{B_j\mid S_j})^2}{d_v}.
\end{equation*}
Apply \cref{eq:group-diagonal-temporal-telescope} with $k+1$ in place of
$k$.  The oracle and search claim is then exactly
\cref{cor:additive-group-locator}.
\end{proof}

\subsection{Square-root wake-ups from the two telescopes}
\label{sec:square-root-wakeups}

The diagonal-loss and correction-energy ledgers should not be charged
separately.  Their product is the actual uncertainty radius of a delayed
response.  This gives a square-root scheduling law before any particular
dynamic Schur representation is chosen.

\begin{lemma}[Two-ledger square-root wake-up bound]
\label{lem:two-ledger-square-root-wakeup}
Use the exact nested corrections and normalized frontier bases of
\cref{thm:causal-orthogonal-response-sketches}.  Fix a vertex $v$ throughout
an interval of events on which it remains exterior.  Partition those events
into disjoint consecutive epochs $I_1,\ldots,I_K$, and put
\begin{equation}
 A_r(v):=\sum_{j\in I_r}\lambda_{jv}^2,
 \qquad
 E_r:=\sum_{j\in I_r}\norm{\Delta^j}_Q^2.
 \label{eq:wakeup-two-ledgers}
\end{equation}
Then the normalized exterior-demand change accumulated at $v$ in epoch $r$
satisfies
\begin{equation}
 \frac{1}{\sqrt{d_v}}
 \left|e_v^TQ\sum_{j\in I_r}\Delta^j\right|
 \leq\sqrt{A_r(v)E_r}.
 \label{eq:wakeup-product-interval}
\end{equation}
Suppose epoch $r$ ends with a certified wake-up only after this product
radius reaches a declared positive margin $\mu_r$, so that
$\sqrt{A_r(v)E_r}\geq\mu_r$.  If
\begin{equation*}
 \mathcal E:=\sum_j\norm{\Delta^j}_Q^2,
 \qquad
 q_{\rm off}:=\frac{1-\alpha}{2},
\end{equation*}
then
\begin{equation}
 \sum_{r=1}^K\mu_r
 \leq
 \sqrt{\frac{q_{\rm off}}{d_v}\,\mathcal E}.
 \label{eq:wakeup-margin-budget}
\end{equation}
In particular, if every completed epoch uses margin at least $\mu>0$, then
\begin{equation}
 K\leq
 \frac{1}{\mu}
 \sqrt{\frac{q_{\rm off}}{d_v}\,\mathcal E}.
 \label{eq:wakeup-count}
\end{equation}

For a unit-mass seed and a trace starting from the zero face, the RPPR
objective gives $\mathcal E\leq\alpha$.  Hence at the finite-band scale
$\mu=c\alpha\tau$ for any fixed $c>0$,
\begin{equation}
 K=O\!\left(\frac{1}{\tau\sqrt{\alpha d_v}}\right).
 \label{eq:wakeup-product-scale}
\end{equation}
\end{lemma}

\begin{proof}
For each $j$, write $\Delta^j=U_ja_j$.  The normalized bases are mutually
$Q$-orthonormal, so
$\norm{a_j}_2=\norm{\Delta^j}_Q$.  Stack the row vectors
$e_v^TQU_j/\sqrt{d_v}$ and the coefficient vectors $a_j$ over
$j\in I_r$.  Ordinary Cauchy--Schwarz gives
\cref{eq:wakeup-product-interval}.

The epochs are disjoint.  A second Cauchy--Schwarz inequality and the wake-up
rule give
\begin{equation*}
 \sum_{r=1}^K\mu_r
 \leq\sum_{r=1}^K\sqrt{A_r(v)E_r}
 \leq
 \sqrt{\left(\sum_rA_r(v)\right)
             \left(\sum_rE_r\right)}.
\end{equation*}
By \cref{thm:source-leverage-diagonal-loss}, the first factor is the
normalized terminal-Schur diagonal loss over a subinterval of the exterior
lifetime of $v$ and is at most $q_{\rm off}/d_v$.  The second factor is at
most $\mathcal E$.  This proves
\cref{eq:wakeup-margin-budget,eq:wakeup-count}.

It remains to record the unit-seed energy budget.  The unregularized smooth
quadratic has minimum value
$-\tfrac12 b^TQ^{-1}b$.  Since
$Q^{-1}\preceq\alpha^{-1}I$ and
$b=\alpha D^{-1/2}s$ with $s\geq0$, $\one^Ts=1$, and $d_v\geq1$,
\begin{equation*}
 b^TQ^{-1}b
 \leq\alpha^{-1}\norm{b}_2^2
 \leq\alpha.
\end{equation*}
The nonnegative $\ell_1$ term can only raise the RPPR objective.  Therefore
the total objective drop from zero is at most $\alpha/2$, while
\cref{eq:orthogonal-correction-telescope} identifies twice that drop with
$\mathcal E$.  Thus $\mathcal E\leq\alpha$, and substituting
$\mu=c\alpha\tau$ proves \cref{eq:wakeup-product-scale}.
\end{proof}

The same argument applies to a fixed hierarchy node $C$ that remains
exterior: replace $A_r(v)$ by
$A_r(C):=\sum_{v\in C}\sum_{j\in I_r}\lambda_{jv}^2$ and the scalar response
by its $D_C^{-1/2}$-normalized vector.  The total first ledger is exactly the
aggregate Schur-diagonal loss in
\cref{eq:cumulative-sketch-diagonal-loss}.  Thus a hierarchy should not be
updated after every batch.  Each node should sleep while both its accumulated
diagonal loss and its accumulated orthogonal energy remain small, and wake
only when their product reaches the node's current pruning margin.

This lemma is a genuine square-root gain, but not yet the desired global work
theorem.  Summing the per-label bound naively can still pay for too many live
boundary labels, and maintaining the diagonal-loss ledger itself is the open
operation.  The useful conclusion is narrower: an implementation that
publishes every additive $\alpha\tau$ change discards an available
Cauchy--Schwarz gain.  The correct scheduler is multiplicative in response
leverage and correction energy.

\subsection{Grounded-Laplacian route to an operator heavy-hitter locator}
\label{sec:grounded-laplacian-locator}

The remaining dynamic object is compatible with spectral vertex-sparsifier
machinery after a simple exact transformation.

\begin{proposition}[Grounded-Laplacian representation]
\label{prop:grounded-laplacian-representation}
Put $c_\alpha:=(1-\alpha)/2$ and
\begin{equation}
 H:=D^{1/2}QD^{1/2}=c_\alpha(D-A)+\alpha D.
 \label{eq:grounded-pagerank-matrix}
\end{equation}
Construct an augmented weighted graph on $V\cup\{\bot\}$ by giving every
original edge conductance $c_\alpha$ and adding a ground edge
$(v,\bot)$ of conductance $\alpha d_v$.  After pinning $\bot$, its Laplacian
principal minor is exactly $H$.

For every eliminated set $S$ and surviving set $W:=V\setminus S$,
\begin{equation}
 Q/S=D_W^{-1/2}(H/S)D_W^{-1/2}.
 \label{eq:schur-congruence-grounded}
\end{equation}
Consequently, for $A\subsetneq T$ and $v\notin T$,
\begin{equation}
 (Q/A)_{vv}-(Q/T)_{vv}
 =\frac{(H/A)_{vv}-(H/T)_{vv}}{d_v}.
 \label{eq:diagonal-loss-grounded}
\end{equation}
If a normalized-coordinate correction is written as
$\Delta_x=D_T^{1/2}\Delta_z$, then its normalized exterior response is
\begin{equation}
 \frac{(Q\Delta_x)_v}{\sqrt{d_v}}
 =\frac{(H\Delta_z)_v}{d_v},
 \qquad v\notin T.
 \label{eq:boundary-response-grounded}
\end{equation}
\end{proposition}

\begin{proof}
The first identity follows by direct multiplication.  The augmented graph
contributes $c_\alpha L$ on original edges and the diagonal
$\alpha D$ through its ground edges.  Diagonal congruence commutes with
principal Schur complementation: substituting
$Q_{UV}=D_U^{-1/2}H_{UV}D_V^{-1/2}$ into the block formula proves
\cref{eq:schur-congruence-grounded}.  Taking a diagonal difference proves
\cref{eq:diagonal-loss-grounded}.  Finally,
$Q\Delta_x=D^{-1/2}H\Delta_z$; divide its $v$th coordinate by
$\sqrt{d_v}$ to obtain \cref{eq:boundary-response-grounded}.
\end{proof}

\begin{proposition}[Exposed-incidence sufficiency]
\label{prop:exposed-incidence-sufficiency}
Fix a face $T$ and let
$W:=\{v\notin T:Q_{vT}\neq0\}$ be its graph boundary.  For every source set
$B\subseteq T$, all of the following objects are determined by
$Q_{TT}$, $Q_{WT}$, and $D_W$:
\begin{enumerate}[label=(\roman*)]
\item the degree-normalized response matrix
\begin{equation}
 \mathcal M_{W\leftarrow B}^{(T)}
 :=D_W^{-1/2}Q_{WT}Q_{TT}^{-1}E_B
   \left(E_B^TQ_{TT}^{-1}E_B\right)^{-1/2};
 \label{eq:exposed-normalized-response}
\end{equation}
\item every source-aware leverage and group mass formed from its row norms;
\item the normalized exterior demand caused by any dual source supported on
      $B$; and
\item every diagonal loss across a nested elimination $A\subsetneq T$, for
      labels in $W$.
\end{enumerate}
Consequently, changing arbitrary edges with both endpoints in $V\setminus T$
while keeping the degrees of labels in $W$ fixed changes none of these
objects.

In the adjacency-access model, scanning the rows of $T$ reveals every
internal edge, cut incidence, and boundary label.  One degree query for each
distinct label in $W$ completes the record.  Its size is
\begin{equation}
 O\!\left(|T|+\sum_{u\in T}d_u+|W|\right)
 =O\!\left(\operatorname{cvol}(T)\right).
 \label{eq:exposed-record-size}
\end{equation}
If a static operator-locator build on a record of size $M$ costs
$\widetilde O(MF)$ for a parameter factor $F$, rebuilding it only at anchor
faces whose charged volumes grow by a fixed constant factor costs
$\widetilde O(VF)$ through final exposed volume $V$.
\end{proposition}

\begin{proof}
The matrix in \cref{eq:exposed-normalized-response} is written entirely in
terms of the three declared blocks.  Its row norms give (ii), and multiplying
the same response by an arbitrary coefficient vector gives (iii).  For (iv),
\cref{eq:source-leverage-diagonal-loss} identifies every required diagonal
loss with the corresponding source-aware squared leverage.  Equivalently, it
is a diagonal of the response product in the block Schur formula.  No block
with both indices in $V\setminus T$ appears.

For the shared unweighted graph, $Q_{TT}$ is fixed by the degrees in $T$ and
the edges internal to $T$, while $Q_{WT}$ is fixed by the cut incidences and
the endpoint degrees.  Scanning the adjacency lists of $T$ reveals the first
two edge sets and all labels in $W$.  Moreover $|W|$ is at most the number of
cut incidences, hence at most $\sum_{u\in T}d_u$, proving
\cref{eq:exposed-record-size}.  The final claim is the geometric-series
argument of \cref{lem:geometric-rebuild} applied to the successive record
sizes.
\end{proof}

\begin{lemma}[Transposed fixed-anchor response sketch]
\label{lem:transposed-anchor-response-sketch}
Fix an anchor $A$, put $U:=V\setminus A$, and let
$\Pi\in\mathbb R^{p\times U}$ be any linear measurement matrix.  Define the
precomputed harmonic sketch
\begin{equation}
 Z_\Pi
 :=Q_{AA}^{-1}Q_{AU}D_U^{-1/2}\Pi^T
 \in\mathbb R^{A\times p}.
 \label{eq:transposed-harmonic-sketch}
\end{equation}
For any frontier $F\subseteq U$ and vector $y\in\mathbb R^F$, lift $y$ by
\begin{equation*}
 \Delta_A:=-Q_{AA}^{-1}Q_{AF}y,
 \qquad
 \Delta_F:=y,
 \qquad
 \Delta_{U\setminus F}:=0,
\end{equation*}
and put $q:=Q\Delta$.  Then $q|_A=0$ and
\begin{equation}
 \Pi D_U^{-1/2}q_U
 =\Pi D_U^{-1/2}Q_{UF}y-Z_\Pi^TQ_{AF}y.
 \label{eq:transposed-response-sketch}
\end{equation}
If $y$ is the exact Schur-frontier solution
\begin{equation*}
 \left(Q_{FF}-Q_{FA}Q_{AA}^{-1}Q_{AF}\right)y=h_F,
\end{equation*}
then $q_F=h_F$.  With $R:=U\setminus F$, the exterior-only measurement is
therefore
\begin{equation}
 \Pi_R D_R^{-1/2}q_R
 =\Pi D_U^{-1/2}q_U-\Pi_FD_F^{-1/2}h_F.
 \label{eq:frontier-contamination-subtraction}
\end{equation}

The two products $Q_{UF}y$ and $Q_{AF}y$ in
\cref{eq:transposed-response-sketch} are generated by scanning the rows of
$F$.  Thus, after the $p$ anchor right-hand sides in
\cref{eq:transposed-harmonic-sketch} have been solved and represented, the
identity requires no dense read of $\Delta_A$ and no pass over $U$.  It does
not declare the multiplication by the touched rows of $Z_\Pi$ free: their
storage, reads, sketch dimension $p$, and every anchor solve remain explicit
response charges.
\end{lemma}

\begin{proof}
The old-face block vanishes by construction.  On $U$,
\begin{equation*}
 q_U=Q_{UF}y-Q_{UA}Q_{AA}^{-1}Q_{AF}y.
\end{equation*}
Left-multiply by $\Pi D_U^{-1/2}$, use symmetry of $Q$, and transpose the
definition of $Z_\Pi$ to obtain
\cref{eq:transposed-response-sketch}.  On $F$, the same block multiplication
gives
\begin{equation*}
 q_F=\left(Q_{FF}-Q_{FA}Q_{AA}^{-1}Q_{AF}\right)y=h_F.
\end{equation*}
The measurement over $U$ is the sum of its $F$ and $R$ column restrictions,
so subtracting the known frontier contribution proves
\cref{eq:frontier-contamination-subtraction}.  Finally, every nonzero of
$Q_{VF}y$ is exposed by a row scan of $F$; the stated implementation ledger
follows directly from the two right-hand-side products.
\end{proof}

This embedding identifies concrete prior machinery, but does not silently
import its runtime.  Van den Brand, Gao, Jambulapati, Lee, Liu, Peng, and Sidford,
``Faster Maxflow via Improved Dynamic Spectral Vertex Sparsifiers,''
arXiv:2112.00722v1, construct a dynamic spectral Schur complement in
Section~4, especially Theorems~4.1 and~4.10, and an operator $\ell_2$
heavy-hitter locator for electrical-flow coordinates in Section~5.4,
Theorem~5.10 and Algorithm~5.  Its blueprint is directly relevant:

\begin{enumerate}
\item precompute sparse heavy-hitter projections of a harmonic-response
      operator;
\item maintain those projected responses through a dynamic terminal Schur
      complement;
\item recover a small candidate set and evaluate only those coordinates.
\end{enumerate}

The source theorem is not the missing local PageRank theorem.  It starts from
the full $m$-edge graph, has initialization
$\widetilde O(m\beta^{-2}\xi^{-2})$, update
$\widetilde O(\beta^{-2}\xi^{-2})$, and locator-query
$\widetilde O(\beta m\xi^{-2})$ in its notation, permits only
$O(\beta m)$ updates before rebuilding, reports edge coordinates relative to
a global $\ell_2$ norm, and initially assumes an oblivious adversary.  The
adaptive reduction in its Section~6 is substantial.  None of these charges is
automatically $\widetilde O(V/\sqrt\alpha)$ or
$\widetilde O(V)$, and \cref{eq:boundary-response-grounded} is a
degree-scaled vertex divergence rather than one edge current.

Nevertheless, this is now the most concrete construction route for the open
oracle.  The new task is not to invent a generic hierarchy from scratch.
\Cref{prop:exposed-incidence-sufficiency} already replaces full-graph input by
an exposure-sized record and makes static geometric rebuilding local.  What
remains is to specialize the operator-heavy-hitter architecture to the
grounded PageRank graph, obtain exposure-charged updates and queries within an
epoch, use \cref{lem:two-ledger-square-root-wakeup} to schedule those queries,
and evaluate only finite-band vertex candidates.  Failure at one of these
named steps becomes an explicit obstruction rather than a vague ``dynamic
SDD'' difficulty.

\Cref{lem:transposed-anchor-response-sketch} gives the corresponding concrete
fixed-anchor interface.  Candidate CountSketch, sparse-recovery, or
hierarchical measurement rows may be placed in $\Pi$.  A geometric rebuild
precomputes their harmonic extensions on the anchor; every aggregate frontier
solve then updates the measured dual signature without forming the old-face
primal correction.  The unresolved quantitative question is whether a
polylogarithmic or output-charged measurement family suffices and whether its
harmonic responses can be updated through an epoch within the target ledger.
The lemma proves the algebra needed for such a construction, not those two
complexity assertions.

The causal-sketch theorem removes three possible sources of repeated work: old numerical
faces are never re-optimized, past random source probes are never replayed,
and a newly exposed boundary label needs no retrospective initialization.
There are now two equivalent remaining operations.  One may apply the $r_j$
signed normalized responses in
\cref{eq:append-only-group-mass-increment} and add their squared mass to the
live hierarchy, or directly maintain the monotone aggregate Schur-diagonal
loss in \cref{eq:cumulative-sketch-diagonal-loss}.  The second formulation
removes the random source dimension and permits scale-matched additive error.
Proving a local bound for either dynamic operation without visiting every leaf
is still open.

The sketch oracle is Monte Carlo.  It gives the displayed simultaneous
one-sided upper estimates only on the probability-$1-\delta$ event.  A
deterministic theorem, or a Las Vegas implementation, additionally needs a
charged validation mechanism for every discarded hierarchy node; no such
general cyclic-core mechanism is claimed here.

\begin{proposition}[Uniform energy intervals do not locate the changed row]
\label{prop:uniform-interval-obstruction}
For every $m\geq1$, there is a PageRank Stieltjes matrix on $m$ disjoint
edges, a fixed active side $S=\{u_1,\ldots,u_m\}$, and exterior side
$W=\{v_1,\ldots,v_m\}$ with the following property.  There are $m$
singleton-source corrections of equal energy such that correction $k$ changes
only exterior coordinate $v_k$.  Nevertheless, the graph-universal or exact
source-oblivious leverage interval marks every not-yet-changed exterior
coordinate ambiguous at every round when all base demands lie half an update
below the reporting threshold.  Refreshing all such ambiguous coordinates
costs $\Omega(m^2)$ queries, whereas the source-aware leverages identify one
possible changed row per round.
\end{proposition}

\begin{proof}
Put $q_0=(1+\alpha)/2$ and $q_1=(1-\alpha)/2$.  On each edge
$u_i v_i$, the shared normalized PageRank block is
\begin{equation*}
 \begin{bmatrix}q_0&-q_1\\-q_1&q_0\end{bmatrix}.
\end{equation*}
Hence $Q_{SS}=q_0I$ and $Q_{WS}=-q_1I$.  For a fixed $h>0$, take correction
$\Delta^k=he_k$.  Its energy radius is $h\sqrt{q_0}$, while its exterior
demand change is $q_1he_k$.  Every source-oblivious leverage equals
$q_1/\sqrt{q_0}$, so the corresponding interval half-width is $q_1h$ at
every exterior coordinate.  If every unchanged demand equals
$\vartheta-q_1h/2$ for a threshold $\vartheta$, all $m-k+1$ unchanged
coordinates have intervals intersecting $\vartheta$ at round $k$ even though
only $v_k$ moves.  Their complete refresh cost sums to
$m+(m-1)+\cdots+1=\Omega(m^2)$.

For the declared source $B_k=\{u_k\}$,
\cref{eq:source-aware-leverage} gives
$\chi_{v_k}^{B_k}=q_1/\sqrt{q_0}$ and
$\chi_{v_j}^{B_k}=0$ for $j\neq k$.  Thus the source-aware intervals isolate
the only coordinate that can change.  On any balanced hierarchy, the exact
group masses in \cref{eq:source-aware-group-mass} also locate that coordinate
with $O(\log m)$ group queries.  This fixed-face construction is a
certificate limitation, not a single-seed RPPR lower bound and not a lower
bound against reporters that exploit the source structure.
\end{proof}

\begin{lemma}[Energy packing of exterior response changes]
\label{lem:response-energy-packing}
Let $S\subseteq T\subseteq S^\star$ be positive exact restricted faces,
extend $x[S]$ by zero on $T\setminus S$, and let $W\subseteq V\setminus T$.
For the exterior demand change
\begin{equation*}
 y_W:=\bigl(c_W-Q_{WT}x[T]_T\bigr)
      -\bigl(c_W-Q_{WT}x[S]_T\bigr),
\end{equation*}
one has $y_W\geq0$ and
\begin{equation}
 \|y_W\|_2^2
 \leq\|x[T]-x[S]\|_{Q_{TT}}^2
 =2\bigl(\phi(x[S])-\phi(x[T])\bigr).
 \label{eq:response-energy-packing}
\end{equation}
If $H_\theta:=\{v\in W:y_v/\sqrt{d_v}\geq\theta\}$, then
\begin{equation}
 \vol(H_\theta)
 \leq\frac{2}{\theta^2}
 \bigl(\phi(x[S])-\phi(x[T])\bigr).
 \label{eq:response-volume-packing}
\end{equation}
The same inequality with the sum of objective shocks holds over disjoint
nested response epochs.
\end{lemma}

\begin{proof}
Put $\Delta:=x[T]_T-x[S]_T\geq0$.  Then
$y_W=-Q_{WT}\Delta\geq0$.  Schur complementation and $Q\preceq I$ give
\begin{equation*}
 Q_{WT}Q_{TT}^{-1}Q_{TW}\preceq Q_{WW}\preceq I.
\end{equation*}
Thus $Q_{WT}Q_{TT}^{-1/2}$ is a contraction and
$\|y_W\|_2\leq\|\Delta\|_{Q_{TT}}$.  Completing the square around the exact
face solution $x[T]$ proves the equality.  Each $v\in H_\theta$ contributes
$y_v^2\geq\theta^2d_v$; summing proves
\cref{eq:response-volume-packing}.  Objective drops telescope over disjoint
nested index intervals.
\end{proof}

\subsection{Monotone boundary variation}
\label{sec:boundary-variation}

The reporting problem is not caused by oscillating exact keys.  The next
ledger shows that every still-inactive demand moves in only one direction and
that the total normalized motion is bounded independently of the number of
active-set events.

\begin{lemma}[Monotone demand-variation ledger]
\label{lem:boundary-variation}
Assume the source has unit mass.  Let
$S_0\subsetneq\cdots\subsetneq S_J\subseteq S^\star$ be a trace whose exact
positive restricted solutions $z^j:=x[S_j]$, extended by zero, satisfy
\begin{equation}
 0\leq z^0\leq z^1\leq\cdots\leq z^J\leq x^\star.
 \label{eq:monotone-face-trace}
\end{equation}
Write $c:=b-\alpha\rho D^{1/2}\one$.  For $v\notin S_j$, define its
normalized shifted demand by
\begin{equation}
 g_v^j:=\frac{c_v-(Qz^j)_v}{\sqrt{d_v}}.
 \label{eq:normalized-boundary-demand}
\end{equation}
Then $g_v^{j+1}\geq g_v^j$ whenever $v\notin S_{j+1}$.  Moreover,
\begin{equation}
 \sum_{j=0}^{J-1}\ \sum_{v\notin S_{j+1}}
 d_v\bigl(g_v^{j+1}-g_v^j\bigr)
 \leq\frac{1-\alpha}{2}.
 \label{eq:boundary-total-variation}
\end{equation}
Consequently every inactive exact key crosses any fixed admission threshold
at most once.
\end{lemma}

\begin{proof}
Put $\Delta^j:=z^{j+1}-z^j\geq0$.  If $v\notin S_{j+1}$, then $v$ is
disjoint from the support of $\Delta^j$, so the Stieltjes signs give
\begin{equation}
 g_v^{j+1}-g_v^j
 =-\frac{Q_{vS_{j+1}}\Delta^j_{S_{j+1}}}{\sqrt{d_v}}
 \geq0.
 \label{eq:boundary-monotone-increment}
\end{equation}
For the shared normalized PageRank matrix, every graph edge $uv$ has
$-Q_{uv}=(1-\alpha)/(2\sqrt{d_ud_v})$.  Multiplying
\cref{eq:boundary-monotone-increment} by $d_v$, summing, and then dropping the
restriction on inactive neighbors gives
\begin{align}
 \sum_{v\notin S_{j+1}}d_v(g_v^{j+1}-g_v^j)
 &\leq\frac{1-\alpha}{2}
   \sum_u\frac{\Delta_u^j}{\sqrt{d_u}}
   \sum_{v\in\mathcal N(u)}1 \notag\\
 &=\frac{1-\alpha}{2}\sum_u\sqrt{d_u}\,\Delta_u^j.
 \label{eq:boundary-variation-one-step}
\end{align}
The coordinate increments telescope under
\cref{eq:monotone-face-trace}, hence their sum is at most $x^\star$.
Stieltjes comparison gives $x^\star\leq x^0:=Q^{-1}b$.  Finally,
$QD^{1/2}\one=\alpha D^{1/2}\one$ and
$b=\alpha D^{-1/2}s$ with $\one^Ts=1$ imply
\begin{equation}
 \sum_u\sqrt{d_u}\,x_u^0=1.
 \label{eq:ppr-normalized-mass}
\end{equation}
Summing \cref{eq:boundary-variation-one-step} over $j$ proves
\cref{eq:boundary-total-variation}.  Monotonicity proves the final claim.
\end{proof}

\begin{corollary}[Limit of additive scalar bucketing]
\label{cor:additive-key-ledger}
Suppose an exact-key implementation publishes an update at $v$ only after
the normalized demand accumulated since its previous publication reaches
$\beta>0$, and let $N_v$ be the number of such publications before $v$ is
activated or the trace terminates.  Then
\begin{equation}
 \sum_v d_vN_v\leq\frac{1-\alpha}{2\beta}.
 \label{eq:additive-key-publications}
\end{equation}
With the finite-band scale $\beta=\alpha\tau/4$, this ledger is only
$O(1/(\alpha\tau))$, the classical push scale, rather than the desired
$\widetilde O(1/(\tau\sqrt\alpha))$ scale.
\end{corollary}

\begin{proof}
Charge every publication at $v$ by $d_v\beta$ of previously unpublished
variation and apply \cref{eq:boundary-total-variation}.  At most one final
sub-$\beta$ remainder is stored per discovered boundary label; these distinct
labels are charged to exposed incidence volume because each has an edge to a
scanned active vertex.
\end{proof}

The corollary is an accounting statement, not a lower bound: it assumes that
the dense increments in \cref{eq:boundary-monotone-increment} are already
available.  Its value is to isolate the required compression.  A successful
finite-band hierarchy should treat each Schur correction as an implicit
nonnegative range-add vector, exploit the one-crossing property, and locate
all threshold crossings without publishing additive $\alpha\tau$ quanta or
refreshing the whole old boundary.

\section{Geometric anchors and aggregate repair}
\label{sec:epochs}

Rebuilding after every support addition recreates repeated-prefix work.  A
minimal remedy is to rebuild only after geometric growth.

\begin{lemma}[Geometric rebuild amortization]
\label{lem:geometric-rebuild}
Let $V_j$ be the charged anchor volume after rebuild $j$, with
$V_{j+1}\geq(1+\gamma)V_j$ for a fixed $\gamma>0$.  If rebuild $j$ costs at
most $cV_j$, then all rebuilds through final volume $V$ cost at most
\begin{equation}
 c\sum_jV_j\leq c\frac{1+\gamma}{\gamma}V.
 \label{eq:rebuild-sum}
\end{equation}
\end{lemma}

\begin{proof}
Read the geometrically increasing sequence backward from its final term and
sum the resulting geometric series.
\end{proof}

\begin{lemma}[Probe before rebuild]
\label{lem:probe-before-rebuild}
Fix a settled anchor $A$, frontier $F$, and face $S=A\cup F$.  Suppose a
frontier repair returns an iterate whose declared error certificate is strong
enough to evaluate the terminal and boundary predicates.  Then:
\begin{enumerate}[label=(\roman*)]
\item if the terminal predicate passes, the algorithm may return without
      constructing a response representation for $Q_{SS}^{-1}$;
\item if it fails, constructing the exact response after the probe produces
      the same state $Q_{SS}^{-1}$ as constructing it before the probe;
\item for a fixed support-event trace and fixed heavy-frontier decisions, the
      probe-first schedule performs no more response rebuilds than the
      rebuild-first schedule and adds at most one certified repair per
      scheduled rebuild, plus a possible terminal probe.
\end{enumerate}
\end{lemma}

\begin{proof}
The terminal certificate concerns the returned iterate, not whether an unused
factor was built, proving (i).  The exact response is a function only of the
principal matrix $Q_{SS}$, so its value is independent of every preceding
iterate and repair, proving (ii).  On the fixed event trace, each heavy
frontier that requires further propagation triggers the same rebuild after its
probe; a terminally certified frontier triggers none.  Charging the one probe
to the corresponding rebuild, with one unmatched terminal probe, proves (iii).
\end{proof}

The lemma pays only the rebuild.  Between rebuilds, small face shocks still
need a lazy response representation; explicitly applying every old-face
correction is not covered.

Aggregate repair avoids a second unnecessary repetition.  Within an epoch,
the response mechanism may update support and boundary certificates while the
numerical iterate remains at its anchor value.  The accumulated block demands
form one final right-hand side.  One preconditioned solve at the end of the
epoch repairs their sum.  Linearity makes this final repair identical to the
sum of exact corrections on the fixed terminal face, while avoiding one old-
face iteration pass per event.

\section{Conditional composition theorem}
\label{sec:composition}

For an SPD preconditioner $P_S$ on a current face, define
\begin{equation}
 \kappa_{\rm eff}(S)
 :=\kappa\!\left(P_S^{-1/2}Q_{SS}P_S^{-1/2}\right).
 \label{eq:effective-condition}
\end{equation}

\begin{theorem}[Conditional response-preconditioned work]
\label{thm:conditional-work}
Consider a support-safe nested trace ending at charged volume $V$.  Suppose:

\begin{enumerate}[label=(\roman*)]
\item every adjacency exposure, response update, old-face read, interval
      refinement, rekey, and reported event has total charge
      $U_{\rm resp}(V)$;
\item anchor rebuilds occur only after $(1+\gamma)$-factor charged-volume
      growth and cost $O(V_j)$ at anchor volume $V_j$;
\item one aggregate repair per epoch reaches its declared energy tolerance in
      \[
      \widetilde O\!\left(\sqrt{\kappa_{\rm eff}(S_j)}\right)
      \]
      preconditioned iterations, each charged by $O(V_j)$;
\item \textsc{Finalize} costs $O(V)$ and returns the terminal certificate.
\end{enumerate}

Then the total charged work is
\begin{equation}
 \widetilde O\!\left(
 U_{\rm resp}(V)
 +V\sqrt{\kappa_{\max}}
 +V
 \right),
 \qquad
 \kappa_{\max}:=\max_j\kappa_{\rm eff}(S_j).
 \label{eq:conditional-work}
\end{equation}
\end{theorem}

\begin{proof}
The first and fourth terms are charged by assumptions (i) and (iv).  Apply
\cref{lem:geometric-rebuild} to the rebuild work.  The aggregate repair work is
at most, up to tolerance logarithms,
\[
 \sqrt{\kappa_{\max}}\sum_jO(V_j)
 =O\!\left(V\sqrt{\kappa_{\max}}\right)
\]
by the same geometric series.  Summing the four disjoint ledger categories
proves \cref{eq:conditional-work}.
\end{proof}

The theorem is deliberately an interface theorem.  Assumption (i) contains
the hard graph problem and may not be replaced by an uncharged oracle call.

\subsection{The two endpoints}

With the identity preconditioner, the shared spectral bounds give
$\kappa_{\max}\leq1/\alpha$.  If response work is no larger than the active
scan work, \cref{eq:conditional-work} has the iterative target
\begin{equation}
 \widetilde O(V/\sqrt\alpha).
 \label{eq:iterative-endpoint}
\end{equation}

At the other endpoint, suppose the response state provides a constant- or
polylogarithmic-quality preconditioner and
$U_{\rm resp}(V)=\widetilde O(V)$.  Then the work approaches
\begin{equation}
 \widetilde O(V).
 \label{eq:response-endpoint}
\end{equation}
This conclusion is conditional on the complete reporting interface, not just
on the existence of a static nearly-linear SDD solver.

\section{Proposed controller}
\label{sec:algorithm}

The following pseudocode is the research contract.  The dense reference backend
implements its algebra and switching order, but not the required local response
or reporting data structure.  All thresholds and certificates remain
note-scoped.

\begin{algorithm}[h]
\caption{Response-preconditioned nested continuation}
\label{alg:response-preconditioned}
\begin{algorithmic}[1]
\STATE Initialize a support-safe face $S$, anchor $A\gets S$, and frontier
buffer $F\gets\emptyset$.
\STATE Build the initial response state and certified boundary dictionary.
\WHILE{the boundary reporter certifies a nonempty violating batch $T$}
    \STATE Expose and charge $T$ and its required incidences.
    \STATE Set $S\gets S\cup T$ and $F\gets F\cup T$.
    \STATE Compute one certified Schur-frontier repair and store its lifted
      response contribution; do not restart a solve on an earlier face.
    \STATE Update every affected certified boundary interval.
    \IF{more propagation is needed and
          $\operatorname{cvol}(F)\geq\gamma\operatorname{cvol}(A)$}
        \STATE Rebuild the anchor: $A\gets A\cup F$ and $F\gets\emptyset$.
    \ENDIF
\ENDWHILE
\STATE Perform one final aggregate repair and materialize the terminal vector.
\STATE Return the vector, boundary certificate, and complete work ledger.
\end{algorithmic}
\end{algorithm}

There are four nonnegotiable safeguards.

\begin{enumerate}
\item The reporter must certify completeness outside its declared finite band.
\item A failed structural fast path must fall back without discarding already
      charged exposure or claiming an unproved transfer of state.
\item Boundary decisions after support growth use only a converged repair; the
      one-probe rule may certify an easy heavy frontier before any rebuild.
\item The final output and every adjacency scan are charged explicitly.
\end{enumerate}

\section{Relation to existing project lines}
\label{sec:map}

\begin{center}
\small
\begin{tabularx}{\textwidth}{@{}p{0.25\textwidth}p{0.26\textwidth}X@{}}
\toprule
Project line & Contribution to this interface & Remaining obligation \\
\midrule
Incremental active-set SDD
& Exact block correction, energy telescope, path-linear implicit solver
& General solve-and-boundary state without repeated materialization \\
Propagate--settle
& Settlement absorption, finite-band reduction, aggregate repair
& Cyclic finite-band response hierarchy \\
Delayed-reflection ladder
& Directed tree, block, shell, and quotient responses
& Nonequitable thick-core event maintenance \\
Volume-gated acceleration
& Safe support cap and one-sided expansion gate
& No-restart expanding-subspace acceleration \\
AESP--CD
& Product-scale local work on a fixed certified envelope
& Accelerated continuation with safe centers \\
Evolving-support CG
& Finite propagation and geometric restart ledger
& Avoid unbounded nonviolating halo exposure \\
\bottomrule
\end{tabularx}
\end{center}

The empirical SOR ladders remain useful propagation arms.  Their exact unit
coordinate rungs are not response factorizations.  They enter the present
framework only after an exact or certified block-response operation is added.

\section{Focused proof program}
\label{sec:proof-program}

\subsection{Finite-band cyclic response}

The strongest immediate target is a data structure that reports every
boundary demand whose normalized lower interval exceeds the gate, certifies
all demands below the safe rejection threshold, and refines only the finite
transition band.  Its total charge must include hierarchy construction,
old-face reads, response applications, rekeys, and output.  Coordinate-specific
leverage estimation is no longer part of the target by
\cref{cor:boundary-leverage-one}.

One plausible route is a hierarchical terminal sparsifier or approximate Schur
representation.  Spectral approximation supplies energy control, and
\cref{lem:boundary-interval} converts it to queried coordinate intervals.  The
search combinatorics reduce, by
\cref{lem:source-aware-hierarchical-locator}, to constant-factor group traces
of source-conditioned response rows.  The unresolved step is to maintain
those traces dynamically on arbitrary cyclic cores and then evaluate all
reached leaves.  By \cref{thm:source-leverage-diagonal-loss}, the same group
quantity is exactly a weighted aggregate decrease of terminal Schur
diagonals; \cref{cor:additive-group-locator} requires only one-sided additive
accuracy at the current pruning scale.  Thus source-response sketches are an
optional estimator, not a second necessary state.  The reporter's unavoidable
output is already packed:
\cref{lem:response-energy-packing} charges every normalized
$\theta$-movement by degree volume to objective shock, while
\cref{lem:boundary-variation} proves one-sided motion and a global weighted
variation ledger.  Before an exact correction is materialized,
\cref{lem:source-aware-leverage-packing} also packs the rows whose certified
intervals can be wide.  By \cref{cor:associative-dual-batching}, a hierarchy
may merge consecutive light sources without changing the exact trajectory at
its checkpoints.  The remaining primitive is therefore a margin-adaptive
aggregate Schur-diagonal-loss oracle, equivalently a source-aware $\ell_2$
locator, over a dynamically maintained terminal representation, with
initialization and every update charged to exposed local volume.

\subsection{Heavy/light frontier events}

Use a dichotomy at every epoch.  A heavy batch causes geometric volume growth
and pays a rebuild.  Light batches remain in the low-dimensional frontier and
are represented lazily until their aggregate correction justifies one repair.
The proof must charge a long sequence of light events without requiring a
full old-face pass after each event.

\subsection{No-restart accelerated repair}

For the response-preconditioned architecture, the algebraic accumulation
problem is closed by \cref{thm:orthogonal-frontier-lifts}: independently
approximated frontier corrections remain orthogonal after lifting, and their
energy errors add exactly.  The implementation obligation is to apply those
lifts within $U_{\rm lift}$ and to maintain the finite-band reporter.

The broader response-free expanding-subspace question remains open.  If the
face grows during one global accelerated recurrence, preserve a support-safe
lower certificate separately from signed accelerated variables, extend new
coordinates with zero numerical and momentum state, and prove that the
estimate-sequence potential survives the enlarged feasible subspace.  The
lift theorem does not provide that result because an exact lift may itself be
a dense old-face response.

\subsection{Lower-bound separation}

Any negative theorem should name the forbidden response resource.  Adjacency
access alone permits the path and tree elimination messages that already beat
repeated-prefix bounds.  A useful separation would prove a product lower bound
for a local-linear-span or bounded-recurrence class and an output-sensitive
response upper bound on the same family.

\section{Experimental contract}
\label{sec:experiments}

A future implementation should compare three backends behind one controller:
iteration with the identity preconditioner, persistent response, and the mixed
method.
Every run should record:

\begin{itemize}
\item newly exposed charged volume and peak active volume;
\item propagation row operations and iterative repair iterations;
\item response updates, factor rebuilds, old-face reads, and rekeys;
\item certified, rejected, and ambiguous boundary intervals;
\item final output writes and verifier work;
\item the exact note-scoped accuracy convention and all solver parameters.
\end{itemize}

Endpoint paths, stars, spiders, tailed fans, high-degree decoys, asymmetric
trees, cycles, equitable thick shells, bounded-block compositions, and
nonequitable cyclic cores exercise different parts of the ledger.  A backend
that is fast only because it omits boundary reporting or final materialization
does not pass the comparison.

\subsection{Dense reference backend and first measurement}

The project now contains an exact dense response--hybrid diagnostic module.
It stores the principal inverse on the anchor, runs warm-started CG on the
exact frontier Schur complement, and absorbs the frontier when combined degree
volume reaches a chosen rebuild factor.  Rebuild factor one is the
pure-response endpoint; infinite factor is the fixed-anchor iterative
endpoint.  This implementation materializes $Q$ and reads every boundary
coordinate, so it validates the algebra and switching logic but supplies no
local-complexity evidence.

At $\alpha=0.05$, note-scoped $\epsppr=10^{-6}$, and seed 7, the factor-two
variant gave the following separated ledger.  ``Resp.'' is the diagnostic
dense response-update flop estimate; ``CG'' is the number of frontier Schur-CG
iterations.  The two columns are deliberately not added.

\begin{center}
\small
\begin{tabular}{@{}lrrrr@{}}
\toprule
Graph & Resp., every batch & Resp., mixed & CG, never rebuild & CG, mixed \\
\midrule
Path, 511 vertices & 129809 & 30101 & 238 & 112 \\
Star, 510 leaves & 133692421 & 1 & 1 & 1 \\
Binary tree, 511 vertices & 133519701 & 16602965 & 36 & 8 \\
Spider, $8\times64$ & 5728401 & 1803761 & 194 & 90 \\
Random 4-regular, 512 vertices & 121475342 & 41788069 & 61 & 40 \\
\bottomrule
\end{tabular}
\end{center}

The mixed arm gives every new frontier one converged Schur-CG probe before a
volume-triggered rebuild.  The path, tree, spider, and regular cases then show
an interior response/iteration tradeoff.  The star collapses safely to the
iterative endpoint: its one heavy batch certifies in one CG step, so no dense
rebuild occurs.  Thus the architecture makes both resources available, but an
individual execution need not consume both.  A universal immediate-rebuild
factor is not justified by these measurements.

\section{Conclusion}
\label{sec:conclusion}

The common object is the restricted Green operator, and preconditioning is the
mathematical bridge between the project's two solver families.  The mixed
architecture is presently the most robust research plan because it can use
exact response where structure is exposed and iterative repair where exact
transport is too expensive.  It is not proved universally best.  The decisive
next theorem is now specifically dynamic local maintenance of additive
aggregate terminal-Schur diagonal losses, or equivalently a local
implementation of the per-batch squared-response range-add on the live
boundary hierarchy,
yielding a fully charged heavy-change reporter for the
energy-packed, monotone finite-band interval crossings on general cyclic
cores.  Full exterior adjacency is not needed: the response object is fixed by
the exposed face, its cut incidences, and boundary degrees, so geometric
static rebuilding is already chargeable to final exposed volume.  The open
cost is now epoch-internal update and locator query work.  The diagonal losses
telescope per label, combine with correction energy through a square-root
wake-up budget, and need only scale-matched one-sided accuracy; neither search
combinatorics, source conditioning, coordinate leverage, nor oscillating keys
remain an obstacle.  At a fixed anchor, transposed harmonic measurements also
remove dense old-face materialization and boundary scans; choosing an
output-charged measurement family and maintaining its responses are the
remaining quantitative obligations.  Numerical
frontier-error accumulation is also closed by
the orthogonal lift theorem.  Response-free continuation over safely expanding
subspaces without repeated restarts remains a separate iterative-track
question.

\end{document}
